% !TeX root = ../../Book4.tex
% 纲要标题：## 第6章*　Teichmüller 空间与曲线的形变
% 纲要位置：几何学计划纲要.md，第 3679 行。
% 正文状态：初稿；尚未正式审核。

% 纲要内容开始
% 
% > 本章目标与证明范围：将曲线模参数从高维复形变中分离。无穷小空间和实际局部参数的联系给明确证明任务；一般可测 Riemann 映照及全局 Teichmüller 深层结果作为逐项外部输入。
% 
% 纲要内容结束

\OptionalChapter{Teichmüller 空间与曲线的形变}{g4:ch:06}
\BookChapterNumber{g4:ch:06}

\begin{chapterabstract}
本章研究曲面的复结构随参数变化时，怎样区分坐标的改变与几何对象的改变。先用拓扑标记消除重新选择道路的歧义，构造映射类群、Dehn 扭转及其同调作用。随后以 Beltrami 方程描述复结构，以向量场产生的平凡形变构造切空间，并用二次微分和 Serre 对偶计算维数。局部参数与有限自同构商说明无穷小计算如何反映实际形变；双曲裤子的长度与扭转提供几何参数。最后将四点球面的交比、光滑族和三个稳定结点纤维完整写出，连接标记碰撞与长颈退化。
\end{chapterabstract}

\begin{learninggoals}
\begin{itemize}
	\item 区分拓扑标记、点的编号、同位类及双全纯同构，解释映射类群商与稳定子。
	\item 在定向环带中构造 Dehn 扭转，计算它对交叉配对和一阶同调的作用。
	\item 推导 Beltrami 微分的坐标变换，构造光滑小参数解，并说明一般可测解使用的外部定理。
	\item 用向量场识别平凡形变，通过 Dolbeault 上同调与 Serre 配对求出切空间、余切空间和维数。
	\item 用双曲裤子的边界长度及实扭转量完成有限拼接与参数恢复，解释标记对扭转整数的记录。
	\item 由交比构造四点球面族，求三种稳定边界、置换作用及结点附近的颈部模量。
\end{itemize}
\end{learninggoals}

把一个平行四边形的对边粘合，便得到复环面。平行四边形的形状改变时，环面的复结构通常也改变；但换一组格基又可能把同一个环面写成另一组参数。这里同时出现两类变化：复结构本身的变化，以及描述它时所选道路的变化。若把两条基本道路始终保留下来，参数问题会比较清楚；最后再忘掉它们，才得到没有拓扑标记的模空间。

\ProseParagraphBreak

第五章为每张曲面给出了统一化模型。现在固定底下的光滑曲面，比较邻近的复结构。在一张复坐标图中，新的全纯坐标满足 \(w_{\bar z}=\mu w_z\)。系数 \(\mu\) 衡量新旧复方向的差异，却仍受到换坐标的影响。向量场的微小流产生 \(\bar\partial\)-恰当的变化；把它们除去后，剩余的类才是几何形变的无穷小参数。与这些类配对的对象是全纯二次微分。这使形变的维数问题回到第三章已经建立的上同调和 Riemann--Roch 计算。

\ProseParagraphBreak

局部参数可以从解析方程取得，也可以从双曲曲面的切割和粘合看见。沿有限条闭测地线剪开曲面，每块成为三边裤子；记录边长以及重新粘回时的错位量，就得到长度与扭转参数。进一步缩短选定的核心曲线，颈部逐渐加长，极限则允许出现结点。四点球面的模型让这一过程只用一个复参数就能写出全部局部方程。

\ProseParagraphBreak

本章使用前三章的全纯坐标、线丛与 Serre 对偶，第五章的统一化、双曲度量和环域模量；光滑曲面及其同位来自第一卷。光滑小形变的解析构造和有限拼接在相应节展开。映射类群可对照 Farb--Margalit 第2.1节与第13.3节，可测方程可对照 Glutsyuk 作者稿第1.2节的陈述及第3节的证明；本章在使用处区分已证构造与逐项外部输入。\cite[第2.1、13.3节]{FarbMargalit2012Primer}\cite[作者稿第1.2、3节]{Glutsyuk2008Uniformization}

% !TeX root = ../../Book4.tex
% 纲要标题：### 6.1 标记与等价关系
% 纲要位置：几何学计划纲要.md，第 3683 行。
% 正文状态：初稿；尚未正式审核。
% 本次扩充的范围、先修与边界：先修为 I.6 的闭曲面分类、II.2 的一阶同调及闭曲面的同调交叉配对，并使用本节已有的标记、同位和群作用。固定闭连通定向曲面，构造 Dehn 扭转并核对方向和交叉数符号，计算其同调作用，将 Torelli 群定义为作用的核。Nielsen–Thurston 三分法限亏格至少二的精确陈述与代表模型；不承担完整分类证明、映射类群的全部生成关系或 Torelli 群的专门结构理论。

% 纲要内容开始
% 
% > 补充写作范围与外部输入：先修为 I.6 的闭曲面分类、II.2 的一阶同调及闭曲面的同调交叉配对，并使用本节已有的标记、同位和群作用。固定闭连通定向曲面，构造 Dehn 扭转并核对方向和交叉数符号，计算其同调作用，将 Torelli 群定义为作用的核。Nielsen–Thurston 三分法限亏格至少二的精确陈述与代表模型；不承担完整分类证明、映射类群的全部生成关系或 Torelli 群的专门结构理论。
% 
% * 标记曲面、同位和双全纯
% * Teichmüller 空间与模空间的区别
% * 映射类群作用及稳定子
% 
% * Dehn 扭转的环带构造与同位类
% * Dehn 扭转对一阶同调的作用
% * Torelli 群作为同调作用的核
% * Nielsen–Thurston 三分法的陈述与模型
% 
% 
% 纲要内容结束

\section{标记与等价关系}
\label{g4:sec:0601}

第四章把亏格一曲面写成格点商，也看到同一条椭圆曲线可以有不同的周期基。要让周期随参数连续变化，必须知道前后两张曲面上的哪条闭路彼此对应。标记提供这种对应；遗忘标记以后，得到的才是只记录双全纯类型的模空间。本节先把这两种等价关系分清，再用环带中的一次扭转说明：改变闭路的对应关系，可以产生同一曲面的不同标记。

\subsection{标记曲面、同位和双全纯}
\label{g4:subsec:0601:01}

固定一张闭、连通、定向光滑曲面 \(S=S_g\)，其亏格为 \(g\)。它作为比较所有曲面的参考曲面，暂不随参数改变。黎曼曲面的复坐标给出实定向；以下的标记和曲面自微分同胚均保持定向。

\begin{definition}\label{g4:def:0601:marked-surface}
	一张以 \(S\) 为参考的\term[biaojiqumian]{标记曲面}[marked surface]是二元组 \((X,f)\)，其中 \(X\) 是紧连通黎曼曲面，\(f:S\to X\) 是保向微分同胚，称为\term[biaoji]{标记}[marking]。两个标记 \(f_0,f_1:S\to X\) 若可由光滑族 \(f_t:S\to X\)、\(0\le t\le1\) 连接，且每个 \(f_t\) 都是保向微分同胚，就称它们\term[tongwei]{同位}[isotopic]，记作 \(f_0\simeq_{\mathrm{iso}}f_1\)。

	两个标记曲面 \((X,f)\)、\((Y,h)\) 等价，是指存在双全纯映射 \(F:X\to Y\)，使
	\begin{equation}\label{g4:eq:0601:marked-equivalence}
		F\circ f\simeq_{\mathrm{iso}}h.
	\end{equation}
\end{definition}

这里的标记是实光滑映射，双全纯条件则加在比较两个复结构的 \(F\) 上。同位过程中的每张映射都要保持可逆；同伦只要求连续映射族，因而条件较弱。例如任意同位都诱导相同的同调同态，但以后会看到，只给定同调同态还不足以记住同位类。

\ProseParagraphBreak

式\eqref{g4:eq:0601:marked-equivalence}确为等价关系。恒等映射给出自反性；将同位逐时取逆并复合 \(F^{-1}\)，得到对称性；把两个双全纯映射复合，并依次连接相应同位，得到传递性。拼接同位时，可先把时间参数光滑重参数化，使其在端点附近恒定，从而保持联合光滑性。

\ProseParagraphBreak

标记会搬运整张曲面上的拓扑资料。若在 \(S\) 上选定辛基 \(a_1,\ldots,a_g,b_1,\ldots,b_g\)，它在 \(X\) 上给出 \(f_*a_i,f_*b_i\)；若还指定了一条简单闭曲线的同位类，标记也把这一同位类搬到 \(X\)。后者比只搬运一个同调类包含更多信息，这一点是引入完整标记的理由。

\ProseParagraphBreak

稍后比较准共形映射、测地线与光滑标记时，会同时遇到连续同伦和光滑同位。下面列出本章采用的曲面拓扑结果，使这些转换有明确依据。

\begin{proposition}\label{g4:prop:0601:surface-isotopy-interface}
	对紧定向光滑曲面，采用以下同伦与同位结果。
	\begin{enumerate}
		\item 每个保向自同胚都拓扑同位于保向微分同胚。两个保向自同胚同伦时可以同位；当二者光滑时，可将同位取为光滑同位。若曲面有边界，且两映射在边界上恒等、给定同伦相对边界，则同位也可相对边界。
		\item 两条本质简单闭曲线自由同伦，当且仅当它们同位；无向曲线的自由同伦允许反转参数方向。对于端点固定在边界上的简单适当弧，也有相对端点的同伦与同位的对应，其中弧的内部始终位于曲面内部。
	\end{enumerate}
\end{proposition}

这些是本章使用的外部曲面拓扑结果：第一项见 \cite[第1.4节，第41--42页，定理1.12、1.13及其后的光滑同位说明；第2.1节，第44--45页]{FarbMargalit2012Primer}；第二项见同书第34页命题1.10及第35--36页的适当弧版本。对光滑曲线，所得嵌入同位再由第二卷定理15.5.3延拓为环境同位。这里同时保留了闭曲线的自由同伦和适当弧的端点条件，后面切割曲面时将分别使用。

\subsection{Teichmüller 空间与模空间的区别}
\label{g4:subsec:0601:02}

\begin{definition}\label{g4:def:0601:teich-moduli}
	所有以 \(S\) 为参考的标记曲面，按\eqref{g4:eq:0601:marked-equivalence}取等价类，所得集合称为 \(S\) 的\term[Teichmullerkongjian]{Teichmüller 空间}[Teichmüller space]\footnote{Paul Julius Oswald Teichmüller（1913--1943），德国数学家，研究几何函数论、准共形映射与黎曼曲面的模参数。}，记作 \(\mathcal T(S)\) 或 \(\mathcal T_g\)。所有亏格为 \(g\) 的紧黎曼曲面按双全纯同构取等价类，所得集合称为\term[mokongjian]{模空间}[moduli space]，记作 \(\mathcal M_g\)。遗忘标记给出映射
	\[
		\pi:\mathcal T(S)\longrightarrow\mathcal M_g,
		\qquad [X,f]\longmapsto[X].
	\]
\end{definition}

也可把所有复结构放在同一张 \(S\) 上。设 \(\mathcal J(S)\) 为与既定定向相容的光滑复结构集合：每个复图册确定切丛自同态 \(J\)，满足 \(J^2=-I\)，局部复坐标中 \(J\partial_x=\partial_y\)。用这些张量的 \(C^\infty\) 拓扑赋予 \(\mathcal J(S)\) 拓扑。记 \(\operatorname{Diff}_0(S)\) 为与恒等映射光滑同位的保向微分同胚群，则
\begin{equation}\label{g4:eq:0601:complex-structure-quotient}
	\mathcal T(S)=\mathcal J(S)/\operatorname{Diff}_0(S).
\end{equation}
本章在左侧采用右侧的商拓扑。为核对这个识别，将 \((X,f)\) 送到拉回复结构 \(f^*J_X\)。若 \(Ff\simeq_{\mathrm{iso}}h\)，则 \(\psi=h^{-1}Ff\) 属于 \(\operatorname{Diff}_0(S)\)，且 \(f^*J_X=\psi^*(h^*J_Y)\)。反过来，这样的拉回等式使 \(h\psi f^{-1}:X\to Y\) 的微分复线性，局部 Cauchy--Riemann 方程说明它双全纯，并满足标记等价。后面将用有限维参数描述这个商空间的局部拓扑。

\begin{example}\label{g4:exm:0601:marked-tori}
	取 \(S=\mathbb R^2/\mathbb Z^2\)，并对 \(\tau\in\mathbb H\) 定义
	\[
		E_\tau=\mathbb C/(\mathbb Z+\tau\mathbb Z),\qquad
		f_\tau([x,y])=[x+\tau y].
	\]
	正数 \(\operatorname{Im}\tau\) 保证 \(f_\tau\) 保向。标准两圈经标记后的 \(\mathrm{d}z\) 周期为 \(1,\tau\)。若 \((E_\tau,f_\tau)\) 与 \((E_\sigma,f_\sigma)\) 等价，设比较它们的双全纯映射为 \(F\)。由全纯微分空间的一维性，\(F^*\mathrm{d}w=a\,\mathrm{d}z\)，其中 \(a\ne0\)。同位使对应周期相等，故
	\[
		a=1,\qquad a\tau=\sigma.
	\]
	于是 \(\tau=\sigma\)。然而 \(\mathbb Z+\tau\mathbb Z=\mathbb Z+(\tau+1)\mathbb Z\)，所以 \(E_\tau\) 与 \(E_{\tau+1}\) 是同一格点商。它们在 \(\mathcal M_1\) 中是同一点，在 \(\mathcal T(S)\) 中的上述标记类却不同：第二条标记闭路增加了第一条闭路。
\end{example}

因此，一个参数变化如果同时搬运了闭路，就自然落在 \(\mathcal T(S)\) 中；比较所得曲面的双全纯类型时，再遗忘这种搬运资料。下一小节给出这一步遗忘的群作用解释。

\subsection{映射类群作用及稳定子}
\label{g4:subsec:0601:03}

记 \(\operatorname{Diff}^+(S)\) 为全部保向微分同胚组成的群。若 \(\psi_t\) 是从恒等映射到 \(\psi\) 的同位，则 \(\varphi\psi_t\varphi^{-1}\) 是从恒等映射到 \(\varphi\psi\varphi^{-1}\) 的同位，所以 \(\operatorname{Diff}_0(S)\) 是正规子群。

\begin{definition}\label{g4:def:0601:mapping-class-group}
	商群
	\[
		\operatorname{Mod}(S)=
		\operatorname{Diff}^+(S)/\operatorname{Diff}_0(S)
	\]
	称为 \(S\) 的\term[yingsheleiqun]{映射类群}[mapping class group]。其元素是保向微分同胚的同位类，乘法由映射复合给出。
\end{definition}

以下固定左作用约定
\begin{equation}\label{g4:eq:0601:mcg-action}
	[\varphi]\cdot[X,f]=[X,f\circ\varphi^{-1}].
\end{equation}
同位的 \(\varphi\) 给出同位的标记；若更换 \((X,f)\) 的等价代表，将原比较映射与 \(\varphi^{-1}\) 复合后仍满足\eqref{g4:eq:0601:marked-equivalence}。所以作用良定义。逆元保证复合顺序正确：先作用 \([\psi]\)，再作用 \([\varphi]\)，所得标记为 \(f\psi^{-1}\varphi^{-1}=f(\varphi\psi)^{-1}\)。

\begin{proposition}\label{g4:prop:0601:moduli-quotient}
	遗忘映射的纤维恰是映射类群的轨道，因而作为集合有
	\[
		\mathcal M_g=\mathcal T(S)/\operatorname{Mod}(S).
	\]
	对 \([X,f]\in\mathcal T(S)\)，记 \(\operatorname{Aut}_0(X)\) 为与恒等映射光滑同位的全纯自同构组成的子群。则有自然同构
	\begin{equation}\label{g4:eq:0601:stabilizer}
		\operatorname{Stab}_{\operatorname{Mod}(S)}[X,f]
		\cong\operatorname{Aut}(X)/\operatorname{Aut}_0(X).
	\end{equation}
\end{proposition}
\begin{proof}
	同一轨道上的点保留同一张 \(X\)，故遗忘映射在轨道上常值。反过来，若 \(F:X\to Y\) 双全纯，给定标记 \(f:S\to X\)、\(h:S\to Y\)，取 \(\varphi=h^{-1}Ff\)，则 \(Ff\varphi^{-1}=h\)。因此 \([X,f]\) 经 \([\varphi]\) 作用后等于 \([Y,h]\)。

	对 \(A\in\operatorname{Aut}(X)\)，定义
	\[
		\iota_f(A)=[f^{-1}Af]\in\operatorname{Mod}(S).
	\]
	它是群同态。令 \(\varphi=f^{-1}Af\)，则 \(Af\varphi^{-1}=f\)，故它的像包含在稳定子中。若 \([\varphi]\) 稳定 \([X,f]\)，标记等价给出某个 \(A\in\operatorname{Aut}(X)\)，满足 \(Af\varphi^{-1}\simeq_{\mathrm{iso}}f\)；于是 \([f^{-1}Af]=[\varphi]\)，所以像就是稳定子。核由 \(f^{-1}Af\simeq_{\mathrm{iso}}\operatorname{id}_S\) 的自同构组成，等价地 \(A\simeq_{\mathrm{iso}}\operatorname{id}_X\)，即 \(\operatorname{Aut}_0(X)\)。第一同构定理给出\eqref{g4:eq:0601:stabilizer}。
\end{proof}

\begin{proposition}\label{g4:prop:0601:automorphism-faithful}
	设 \(X\) 是亏格 \(g\ge2\) 的紧连通黎曼曲面。若全纯自同构 \(A:X\to X\) 与恒等映射同伦，则 \(A=\operatorname{id}_X\)。特别地，\(\iota_f\) 单射，且 \([X,f]\) 的稳定子同构于 \(\operatorname{Aut}(X)\)。
\end{proposition}
\begin{proof}
	假设 \(A\ne\operatorname{id}_X\)。恒等定理说明它的不动点孤立；由于 \(X\) 紧，不动点只有有限多个。在不动点 \(p\) 的正向复坐标中写 \(A(z)=a(z)\)、\(p=0\)。全纯函数 \(z-a(z)\) 不恒为零，其零点阶数 \(m_p\) 为正整数。沿足够小的正向圆周，\(z-a(z)\) 的绕数为 \(m_p\)，故实二维局部不动点指数为 \(m_p>0\)。

	第二卷定理16.4.4的 Lefschetz 公式\footnote{Solomon Lefschetz（1884--1972），俄国出生、后在美国工作的数学家，研究代数拓扑与代数几何，建立同调方法的不动点公式。}适用于这张闭定向光滑曲面，给出
	\[
		L(A)=\sum_{p\in\operatorname{Fix}(A)}m_p\ge0.
	\]
	另一方面，同伦于恒等映射使 \(A\) 在各阶实上同调上的作用均为恒等，因此
	\[
		L(A)=1-\dim H^1(X;\mathbb R)+1=2-2g<0,
	\]
	矛盾。这里即使没有不动点，前一个和也只是零，仍然产生矛盾。由同位蕴含同伦，\(\operatorname{Aut}_0(X)\) 平凡，再用\eqref{g4:eq:0601:stabilizer}即可。
\end{proof}

亏格一恰好显示了假设的作用。在 \(E=\mathbb C/\Lambda\) 上，每个平移 \(T_b([z])=[z+b]\) 都是全纯自同构，且 \(T_{tb}\)、\(0\le t\le1\)，给出到恒等映射的同位，因此全纯自同构到映射类群的同态有非平凡核。事实上，这个核恰为平移群。若 \(A\) 同伦于恒等映射，写 \(A^*\mathrm{d}z=a\,\mathrm{d}z\)，沿任一非零格点周期积分得到 \(a\lambda=\lambda\)，所以 \(a=1\)。于是全纯映射 \(P\mapsto A(P)-P\) 的微分处处为零，在连通的 \(E\) 上为常值，故 \(A\) 是平移。这样，环面标记的稳定子对应于全纯自同构群除去平移后的商群。

\subsection{Dehn 扭转的环带构造与同位类}
\label{g4:subsec:0601:04}

在 \(S\) 上取一条光滑简单闭曲线 \(c\)。它有一条环带邻域：第二卷定理15.3.5给出法线丛的管状模型，而 \(S\) 的定向和 \(c\) 的暂定方向使实法线丛可定向，因而平凡。取坐标
\[
	\iota:(\mathbb R/\mathbb Z)\times[-1,1]\longrightarrow S,
	\qquad c=\iota((\mathbb R/\mathbb Z)\times\{0\}),
\]
使 \((\partial_\theta,\partial_t)\) 按此顺序给出 \(S\) 的正定向，并使 \(c\) 沿 \(\theta\) 增大方向定向。

\begin{definition}\label{g4:def:0601:dehn-twist}
	选光滑非减函数 \(\chi:[-1,1]\to[0,1]\)，在 \(-1\) 附近恒为 \(0\)，在 \(1\) 附近恒为 \(1\)。环带上的映射
	\begin{equation}\label{g4:eq:0601:dehn-annulus}
		T_c(\theta,t)=(\theta+\chi(t),t)
	\end{equation}
	在两条边界附近均为恒等，其中 \(\theta\) 按模 \(1\) 计算。把它在环带外延拓为恒等，所得保向微分同胚称为沿 \(c\) 的正\term[Dehnniuzhuan]{Dehn 扭转}[Dehn twist]\footnote{Max Wilhelm Dehn（1878--1952），德国出生、后在美国工作的数学家，研究几何、拓扑与群的判定问题。}。它的逆映射把 \(+\chi(t)\) 换成 \(-\chi(t)\)。
\end{definition}

Jacobi 矩阵为 \(\bigl(\begin{smallmatrix}1&\chi'(t)\\0&1\end{smallmatrix}\bigr)\)，行列式为 \(1\)，因此保向和可逆性都可直接核验。在环带的带状覆盖 \(\mathbb R\times[-1,1]\) 中，一条从下边到上边的竖弧经 \(T_c\) 后，上端比原来多走一个 \(\theta\) 周期；下降到环带，两端仍是原来的点，中间的弧则绕核心多转一周。\cref{g4:fig:0601:dehn-strip}把这一局部运动画在带状覆盖中。

% !TeX root = ../Book4.tex
\begin{figure}[htbp]
	\centering
	\begin{tikzpicture}[x=2.1cm,y=1.15cm,font=\small]
		\fill[GeometryBlueLight] (-.2,-1) rectangle (2.2,1);
		\draw[GeometryGray] (-.2,-1)--(2.2,-1);
		\draw[GeometryGray] (-.2,1)--(2.2,1);
		\draw[GeometryGray,dashed] (0,-1)--(0,1);
		\draw[GeometryGray,dashed] (1,-1)--(1,1);
		\draw[GeometryGray,dashed] (2,-1)--(2,1);
		\draw[GeometryBlue,thick,->] (.25,-1)--(.25,.85);
		\draw[GeometryBlue,thick] (.25,.85)--(.25,1);
		\draw[red!70!black,thick,->] (.25,-1)--(.25,-.75) .. controls (.25,-.3) and (1.25,.3) .. (1.25,.75)--(1.25,.9);
		\draw[red!70!black,thick] (1.25,.9)--(1.25,1);
		\draw[GeometryGray,->] (0,0)--(1.8,0) node[above] {\(c:\ \theta\text{ 增大}\)};
		\draw[<->] (.25,1.25)--(1.25,1.25) node[midway,above] {一个周期};
		\node[left] at (-.2,1) {\(t=1\)};
		\node[left] at (-.2,-1) {\(t=-1\)};
		\node[GeometryBlue,left] at (.25,-.55) {\(\gamma\)};
		\node[red!70!black,right] at (1.2,.55) {\(T_c\gamma\)};
		\node[below] at (0,-1) {\(0\)};
		\node[below] at (1,-1) {\(1\)};
		\node[below] at (2,-1) {\(2\)};
	\end{tikzpicture}
	\caption[Dehn 扭转的带状覆盖]{环带的带状覆盖中，蓝色横弧经正扭转变为红色弧。下端提升不变，上端提升沿核心正向增加一个周期；按 \(\theta\sim\theta+1\) 取商后，两条弧端点相同。图中定向为先向右的核心方向、后向上的横向方向。}
	\label{g4:fig:0601:dehn-strip}
\end{figure}


\begin{proposition}\label{g4:prop:0601:twist-isotopy}
	\(T_c\) 的同位类与函数 \(\chi\)、环带宽度、保向管状坐标及 \(c\) 的暂定方向无关，只依赖 \(c\) 的无向同位类与 \(S\) 的定向。对任意保向微分同胚 \(\varphi\)，有
	\begin{equation}\label{g4:eq:0601:twist-conjugacy}
		[\varphi][T_c][\varphi]^{-1}=[T_{\varphi(c)}].
	\end{equation}
\end{proposition}
\begin{proof}
	两个允许的函数 \(\chi_0,\chi_1\) 的凸组合仍在两端附近分别等于 \(0,1\)。把它代入\eqref{g4:eq:0601:dehn-annulus}，便得到全局微分同胚之间的同位。特别地，可将 \(\chi\) 从 \(0\) 变到 \(1\) 的区域压缩到核心附近任意窄的子环带。因此比较不同管状坐标时，可以先把两个扭转的支撑缩到共同的小邻域。

	将法向微分正规化后，第二卷命题15.3.6给出两张管状坐标间固定核心的光滑同位，且可取共同的正宽度。在每一时刻的管状坐标中使用同一个、在边界附近恒定的 \(\chi\)，再在外部补恒等映射，就得到全局扭转之间的光滑同位。法向坐标的正缩放可先沿正函数的凸组合连接；核心的保向重新参数化也可连接到恒等映射：若其提升为 \(h:\mathbb R\to\mathbb R\)，满足 \(h(x+1)=h(x)+1\)、\(h'>0\)，则 \((1-s)x+sh(x)\) 给出这样的同位。这处理了坐标的各种选择。

	若反转 \(c\) 的方向，保持环境定向的配套坐标变换为 \((\theta,t)\mapsto(-\theta,-t)\)。在旧坐标中，新扭转写成 \((\theta-\chi(-t),t)\)，模 \(1\) 后也就是 \((\theta+1-\chi(-t),t)\)。函数 \(1-\chi(-t)\) 仍满足原来的端点条件，故同位类不变。

	最后，\(\varphi\) 搬运环带坐标，直接把 \(T_c\) 共轭成 \(\varphi(c)\) 附近同方向的扭转，这给出\eqref{g4:eq:0601:twist-conjugacy}。若 \(c\) 经光滑嵌入同位移到 \(c'\)，第二卷定理15.5.3将这个紧圆周的同位延拓为环境同位 \(\varphi_t\)。于是 \([\varphi_1]=1\)，共轭公式给出 \([T_c]=[T_{c'}]\)。
\end{proof}

这里的“正”由坐标式\eqref{g4:eq:0601:dehn-annulus}和定向次序共同规定。将环境定向反转，会交换正、负扭转。只反转核心曲线的方向则不改变扭转，因为横向坐标也同时反转了。

\begin{lemma}\label{g4:lem:0601:relative-disk-annulus}
	闭圆盘上逐点固定边界的保向自同胚，相对边界同位于恒等映射。闭环带 \(A=(\mathbb R/\mathbb Z)\times[-1,1]\) 上逐点固定边界的保向自同胚 \(h\)，其相对边界同位类由一个整数 \(n(h)\) 唯一确定：将 \(h\) 提升到带状覆盖，并使提升 \(\widetilde h\) 在下边界恒等，则上边界满足
	\[
		\widetilde h(x,1)=(x+n(h),1).
	\]
	在本节的方向约定下，\(n(T_c)=1\)，所以每个相对边界同位类唯一写为 \([T_c^n]\)。对光滑自微分同胚，上述同位也可取为光滑同位。
\end{lemma}
\begin{proof}
	先把圆盘识别为 \(\{|x|\le1\}\)。若 \(h\) 固定边界，则对 \(0<t\le1\) 定义
	\[
		H_t(x)=
		\begin{cases}
			t\,h(x/t),&|x|\le t,\\
			x,&t\le|x|\le1,
		\end{cases}
		\qquad H_0(x)=x.
	\]
	两式在 \(|x|=t\) 相同，每个 \(H_t\) 都是固定边界的同胚，其逆用 \(h^{-1}\) 的同一公式给出。又有 \(|H_t(x)-x|\le2t\)，故在 \(t=0\) 联合连续。这样得到恒等映射到 \(h\) 的相对边界同位。光滑情形用\cref{g4:prop:0601:surface-isotopy-interface}的相对版本将同位光滑化。

	对环带，第一卷定理5.3.10应用于单连通的带状覆盖，给出指定一点后的唯一提升；固定下边界使这一提升在整条下边界恒等。上边界的提升是整数平移，整数不可能随 \(x\) 或相对边界同位连续改变。两个映射复合时该整数相加，所以 \(n\) 是同态，且\eqref{g4:eq:0601:dehn-annulus}给出 \(n(T_c^m)=m\)。只须证明 \(n(h)=0\) 时 \(h\) 相对边界同位于恒等映射。

	这时 \(\widetilde h\) 在上下边界都恒等，且
	\[
		\widetilde h(x+1,t)=\widetilde h(x,t)+(1,0).
	\]
	后一式来自 \(h\) 在 \(\pi_1(A)\) 上为恒等：边界圈是其生成元，而 \(h\) 逐点固定该圈。带状区域凸，故 \((1-s)\widetilde h+s\operatorname{id}\) 给出保持两条边界的同伦，并与整数平移交换。它下降为 \(h\) 到恒等映射的相对边界同伦，再用\cref{g4:prop:0601:surface-isotopy-interface}得到同位。对一般 \(h\)，将其复合 \(T_c^{-n(h)}\) 即归结到这一情形。这个圆盘与环带的计算也见 \cite[第47--48页，引理2.1；第51--52页，命题2.4]{FarbMargalit2012Primer}。
\end{proof}

\subsection{Dehn 扭转对一阶同调的作用}
\label{g4:subsec:0601:05}

沿用第四章的交数约定：\(u\cdot v\) 在横截交点处先取 \(u\) 的切向量，再取 \(v\) 的切向量，与曲面的正定向比较。特别地，在刚才的环带中，核心 \(c\) 与沿 \(t\) 增大的横弧按这个次序相交为正。第二卷定理16.2.8及推论16.2.9保证该配对只依赖整数同调类。

\begin{proposition}\label{g4:prop:0601:twist-homology}
	对任意 \(x\in H_1(S;\mathbb Z)\)，正 Dehn 扭转满足
	\begin{equation}\label{g4:eq:0601:twist-transvection}
		(T_c)_*x=x+([c]\cdot x)[c]
		=x-(x\cdot[c])[c].
	\end{equation}
	更一般地，对 \(n\in\mathbb Z\)，有
	\begin{equation}\label{g4:eq:0601:twist-power}
		(T_c^n)_*x=x+n([c]\cdot x)[c].
	\end{equation}
\end{proposition}
\begin{proof}
	一阶同调由闭曲线生成，故先对一条有向光滑闭曲线 \(\gamma\) 核对。作小扰动使它横截于 \(c\)，再把扭转支撑的环带缩窄。此时 \(\gamma\) 在支撑环带内的相关部分，是围绕各交点的有限条横穿弧；其余部分不受扭转影响。

	设一条横穿弧沿 \(t\) 增大方向通过核心。在带状覆盖中，原弧与扭转后的弧的下端可取相同提升，上端提升相差 \(+1\)。将扭转后的弧接上原弧的反路，得到的闭路因而代表 \([c]\)。对应交点的 \(c\cdot\gamma\) 符号也是 \(+1\)。若横穿方向相反，上端与下端的角色交换，差类和交点符号同时变成负号。因此每个交点贡献 \(\varepsilon_p[c]\)，将有限条弧相加便得到
	\[
		[T_c(\gamma)]-[\gamma]
		=\sum_{p\in c\cap\gamma}\varepsilon_p[c]
		=([c]\cdot[\gamma])[c].
	\]
	按整数线性延拓即得\eqref{g4:eq:0601:twist-transvection}。由于 \([c]\cdot[c]=0\)，其中加上的算子 \(N(x)=([c]\cdot x)[c]\) 满足 \(N^2=0\)。故 \((I+N)^n=I+nN\) 对正、零、负整数均成立，得到\eqref{g4:eq:0601:twist-power}。
\end{proof}

公式的两种写法把符号说明得很清楚：若把交数写成“被作用的类先出现”，前面的号就是负号。反转 \(c\) 的方向使 \([c]\) 和 \([c]\cdot x\) 同时反号，所以右侧保持不变。

\begin{example}\label{g4:exm:0601:twist-matrix}
	在一个手柄中取 \(a\cdot b=1\)，其余基曲线避开 \(a\) 的支撑环带，则
	\[
		(T_a)_*a=a,\qquad (T_a)_*b=b+a.
	\]
	在有序基 \((a,b)\) 下，按列记录基向量的像，矩阵为
	\[
		\begin{pmatrix}1&1\\0&1\end{pmatrix}.
	\]
	这正是\cref{g4:exm:0601:marked-tori}中将第二个周期从 \(\tau\) 改成 \(\tau+1\) 的同调变换。若曲线 \(c\) 不分离 \(S\)，沿 \(c\) 切开后的曲面连通，取一条连接两条新边界、且两端在粘合时互相对应的简单弧，粘回并圆滑后得到与 \(c\) 相交一次的闭曲线 \(b\)。适当定向使 \([c]\cdot[b]=1\)，故 \((T_c^n)_*[b]=[b]+n[c]\)。由于 \([c]\ne0\) 且 \(H_1(S;\mathbb Z)\) 无挠，所有 \([T_c^n]\) 两两不同，\([T_c]\) 在映射类群中有无限阶。
\end{example}

\subsection{Torelli 群作为同调作用的核}
\label{g4:subsec:0601:06}

保向微分同胚保持局部相交符号，因而保持同调交数。同位映射又诱导相同的同调同态，所以有群同态
\begin{equation}\label{g4:eq:0601:homology-representation}
	\rho:\operatorname{Mod}(S)\longrightarrow
	\operatorname{Sp}(H_1(S;\mathbb Z),\cdot),
	\qquad [\varphi]\longmapsto\varphi_*.
\end{equation}
右侧是保持该反对称非退化整数配对的自同构群。选辛基 \((a_1,\ldots,a_g,b_1,\ldots,b_g)\) 后，交数矩阵为
\[
	J=\begin{pmatrix}0&I_g\\-I_g&0\end{pmatrix},
\]
右侧即 \(\operatorname{Sp}(2g,\mathbb Z)=\{A\in\operatorname{GL}(2g,\mathbb Z):A^{\mathsf T}JA=J\}\)。

\begin{definition}\label{g4:def:0601:torelli}
	同调作用\eqref{g4:eq:0601:homology-representation}的核
	\[
		\mathcal I(S)=\ker\rho
	\]
	称为 \(S\) 的\term[Torelliqun]{Torelli 群}[Torelli group]。它由在 \(H_1(S;\mathbb Z)\) 上恒等作用的映射类组成，是 \(\operatorname{Mod}(S)\) 的正规子群。
\end{definition}

\begin{example}\label{g4:exm:0601:torelli-twists}
	若简单闭曲线 \(c\) 分离 \(S\)，其一侧闭包是一张带边界曲面，边界为 \(c\)。由第一卷定理6.2.10给该曲面作与边界相容的有限三角剖分，所有内部边在定向三角形之和的边界中相消，留下 \(c\)，所以 \([c]=0\)。于是\eqref{g4:eq:0601:twist-transvection}给出 \((T_c)_*=I\)，即
	\[
		[T_c]\in\mathcal I(S).
	\]
	另设 \(c,d\) 不交，分别定向后满足 \([c]=[d]\)。两个扭转在同调上给出同一个变换，所以
	\[
		[T_cT_d^{-1}]\in\mathcal I(S).
	\]
	若 \(c,d\) 均不分离、不同位且它们的并围住一张子曲面，通常称这一对为\term[bianjiequxiandui]{边界曲线对}[bounding pair]，上述复合称为边界曲线对映射。等式 \([c]=[d]\) 来自该子曲面的定向边界为 \(c-d\)。这些计算直接检验了它们属于核的原因。
\end{example}

Torelli 群也说明两种“忘记资料”的层次。先把标记中的每条曲线只记成同调类，相当于允许来自 \(\mathcal I(S)\) 的改变；把整个标记忘掉，则允许全部映射类群。第四章的归一化周期以同调辛基为依据，因此经过 Torelli 群的标记变化，所有基闭路的周期都保持原样。周期能有效记录复结构，而所选闭路的完整同位资料要在 \(\mathcal T(S)\) 中另外保存。

\subsection{Nielsen–Thurston 三分法的陈述与模型}
\label{g4:subsec:0601:07}

Nielsen--Thurston 分类把映射类与曲面上可保留的几何结构联系起来。\footnote{Nielsen，Jakob Nielsen（1890--1959），丹麦数学家，研究曲面映射与群自同构。Thurston，William Paul Thurston（1946--2012），美国数学家，研究低维拓扑、双曲几何与几何结构。}以下设 \(g\ge2\)。一条简单闭曲线若不零伦，就称为本质曲线。\term[duochongquxian]{多重曲线}[multicurve]是有限条本质简单闭曲线的不交并，且任意两条不同位。映射类保持多重曲线，是指它把这些无向曲线的同位类作为一个集合置换，允许交换其中的分支。

\begin{definition}\label{g4:def:0601:nt-types}
	映射类 \(\phi\in\operatorname{Mod}(S)\) 若有有限阶，就称为\term[zhouqixing]{周期型}[periodic]；若保持某个非空多重曲线，就称为\term[keyue]{可约}[reducible]。为定义第三类，先引入曲面上的带测度奇异叶状结构。

	在除去有限集的曲面上，局部坐标中的平行直线 \(y=\text{常数}\) 给出一维叶状结构；横弧带有测度 \(|\mathrm{d}y|\)，并要求不同坐标中的测度相容。允许的奇点为 \(p\) 叉奇点，\(p\ge3\)：从该点伸出 \(p\) 条分离叶，相邻叶扇区按平行叶的标准模型拼合。这类对象称为\term[daiceduqiyiyezhuangjiegou]{带测度奇异叶状结构}[measured singular foliation]。相互横截的两个此类结构在普通点可同时写成水平叶与竖直叶，在共同奇点有相同的叉数。

	若 \(\phi\) 有保向同胚代表 \(h\)，保持这样一对横截的带测度奇异叶状结构 \((\mathcal F_1,\mu_1)\)、\((\mathcal F_2,\mu_2)\)，并存在 \(\lambda>1\)，使对各自的横弧 \(\alpha\) 有
	\begin{equation}\label{g4:eq:0601:pseudo-anosov-measure}
		\mu_1(h\alpha)=\lambda\mu_1(\alpha),\qquad
		\mu_2(h\alpha)=\lambda^{-1}\mu_2(\alpha),
	\end{equation}
	就称 \(\phi\) 为\term[weiAnosov]{伪 Anosov 型}[pseudo-Anosov]\footnote{\OriginalName{Дмитрий Викторович Аносов}（Dmitrii Viktorovich Anosov，1936--2014），俄罗斯数学家，研究动力系统与遍历理论。}，\(\lambda\) 称为其伸长因子。
\end{definition}

叉数的局部解析模型可写得更明确：在复坐标 \(z\) 中，取二次微分 \(z^{p-2}\,\mathrm{d}z^2\)，在避开原点的小扇区上积分其平方根。积分正比于 \(z^{p/2}\)，其实部与虚部的等值线给出相互横截的两族叶，横向测度取相应微分的绝对值。平方根换号只反转坐标，保持叶和测度；沿原点绕行后因此仍能拼合。水平或竖直方向在原点各有 \(p\) 条分离叶。闭曲面没有穿孔，本节的奇点条件采用 \(p\ge3\)。

\ProseParagraphBreak

为了容许这样的奇点模型，伪 Anosov 代表用同胚表述。\cref{g4:prop:0601:surface-isotopy-interface}把闭曲面的拓扑映射类与光滑映射类自然识别，因而定义中的同胚代表仍指定前面定义的同一个映射类。

\begin{theorem}[Nielsen--Thurston 分类]\label{g4:thm:0601:nielsen-thurston}
	设 \(S\) 是亏格 \(g\ge2\) 的闭连通定向曲面。每个保向映射类恰属于下面三类之一：
	\begin{enumerate}
		\item 周期型。
		\item 非周期的可约型，即该类有无限阶并保持一个非空多重曲线。
		\item 伪 Anosov 型。
	\end{enumerate}
\end{theorem}

本定理采用 \cite[第13.3节，第376页，定理13.2]{FarbMargalit2012Primer} 的分类结果，限制到闭曲面 \(g\ge2\) 的情形。这里把第二项写成“非周期的可约型”，是因为按定义，周期型也可能保持多重曲线；例如恒等映射保持每一条曲线。上述三项的互斥性包含了伪 Anosov 类有无限阶且不保持任何非空多重曲线这一结论。

\begin{example}\label{g4:exm:0601:periodic-reducible-models}
	取无重根多项式 \(P\)，次数为 \(2g+2\)。\cref{g4:prop:0204:hyperelliptic-completion,g4:cor:0204:hyperelliptic-genus}构造的紧曲面 \(X:y^2=P(z)\) 亏格为 \(g\)，且
	\[
		\sigma(z,y)=(z,-y)
	\]
	延拓为二阶全纯自同构。它不是恒等映射，由\cref{g4:prop:0601:automorphism-faithful}，其映射类也不是恒等类，因而给出周期型模型。

	若 \(c\) 为不分离的简单闭曲线，\([T_c]\) 保持 \(c\) 的同位类，由\cref{g4:exm:0601:twist-matrix}又有无限阶，所以给出第二类模型。曲面上的其他区域可以完全不动，而环带中的横穿弧不断增加绕行次数。
\end{example}

伪 Anosov 模型需要同时出现伸长、压缩和有限个叉点。下面从一个完全可计算的环面映射构造亏格二的例子；分歧覆盖把环面上没有奇点的平行叶变成四叉叶。

\begin{example}\label{g4:exm:0601:pseudo-anosov-model}
	在 \(T=\mathbb R^2/\mathbb Z^2\) 上取整矩阵
	\[
		A=\begin{pmatrix}2&1\\1&1\end{pmatrix}.
	\]
	它的行列式为 \(1\)，故诱导保向自微分同胚。两个实特征值为
	\[
		\lambda=\frac{3+\sqrt5}{2}>1,
		\qquad\lambda^{-1}=\frac{3-\sqrt5}{2}.
	\]
	选保向实线性坐标 \((u,v)\) 使其表达为 \((u,v)\mapsto(\lambda u,\lambda^{-1}v)\)。平面上的 \(u=\text{常数}\)、\(v=\text{常数}\) 两族平行线及横向测度 \(|\mathrm{d}u|,|\mathrm{d}v|\) 对整格点平移不变，因而下降到环面；映射 \(A\) 对这些测度的作用恰为\eqref{g4:eq:0601:pseudo-anosov-measure}。

	要得到亏格二，取 \(P=[0,0]\)、\(Q=[1/2,0]\)。计算
	\[
		A^3=\begin{pmatrix}13&8\\8&5\end{pmatrix}\equiv I\pmod2,
	\]
	可见 \(B=A^3\) 固定 \(P,Q\)。在两点穿孔的环面 \(T^\circ=T\setminus\{P,Q\}\) 上，选两条手柄闭路 \(a,b\) 及两条穿孔小圈 \(d_P,d_Q\)。切开成多边形再应用第一卷定理5.2.4的 van Kampen 定理，得到展示
	\[
		\pi_1(T^\circ)=
		\langle a,b,d_P,d_Q\mid[a,b]d_Pd_Q=1\rangle.
	\]
	将 \(d_P,d_Q\) 都送到 \(\mathbb Z/2\mathbb Z\) 的 \(1\)，而将 \(a,b\) 任意送到 \(0\) 或 \(1\)，恰有四个群同态。关系式在目标中成为 \(1+1=0\)，故这四种选择均成立。每个同态都满射，由第一卷定理5.5.6的单值重建得到连通二叶覆盖 \(E\to T^\circ\)。在每个穿孔周围的单值为交换两叶，故\cref{g4:lem:0202:punctured-disk-cover}的局部模型为 \(w=z^2\)。在两个末端各补一个点，便得到到 \(T\) 的二次分歧覆盖
	\[
		\pi:\widehat T\longrightarrow T,
	\]
	恰在 \(P,Q\) 上方简单分歧。给 \(T\) 任取一个平坦复结构，覆盖片与这些补点坐标使 \(\pi\) 全纯，且有限叶覆盖与有限个补点圆盘使 \(\widehat T\) 紧。Riemann--Hurwitz 公式给出
	\[
		2g(\widehat T)-2=2\bigl(2g(T)-2\bigr)+2=2,
	\]
	所以 \(g(\widehat T)=2\)。

	\(B\) 固定两个穿孔，因而通过拉回置换上述四个同态；基点的改变只产生内自同构，对交换目标没有影响。四个对象的任意置换的阶整除 \(12\)，因此可以取 \(k=12\)，使 \(B^k\) 保持所选覆盖的单值同态。覆盖重建的等价性于是给出同胚 \(\widehat B:E\to E\)，满足 \(\pi\widehat B=B^k\pi\)。在每个穿孔的上方只有一个末端，\(\widehat B\) 及其逆将该末端送回自身，因此都连续延拓到补点，得到 \(\widehat T\) 的保向同胚。

	最后拉回两族平行叶及其横向测度。在普通覆盖点，它们仍是相互横截的平行叶；在一个分歧点的坐标 \(w=z^2\) 中，每族穿过 \(w=0\) 的直线的逆像有四条射线，故得到四叉奇点。具体地，原来的两个实线性坐标 \((u,v)\) 组成一个保向实线性局部坐标；在该坐标中，二重分歧模型经保向局部同胚化为平方映射，拉回叶的模型正是 \(p=4\) 的模型。于是两个补点都是共同的四叉奇点，其余点没有奇点。等式 \(\pi\widehat B=B^k\pi\) 使两种横向测度分别乘以 \(\lambda^{3k}\) 和 \(\lambda^{-3k}\)。因此 \(\widehat B\) 是伸长因子为 \(\lambda^{3k}>1\) 的伪 Anosov 代表。
\end{example}

三类模型分别展示有限次迭代后返回、沿保留曲线持续扭转，以及在全曲面上保持两族相反伸缩的奇异叶。分类定理说明每个映射类都落入这三种情形；映射类群作用于标记空间时，这些拓扑运动也就成为比较复结构的三种基本动力学行为。

% 纲要条目排为编号小节标题，并进入目录与 PDF 书签；说明段仍保留于纲要注释中。

% !TeX root = ../../Book4.tex
% 纲要标题：### 6.2 Beltrami 方程与复结构
% 纲要位置：几何学计划纲要.md，第 3697 行。
% 正文状态：初稿；尚未正式审核。

% 纲要内容开始
% 
% * Beltrami 微分和局部变形
% * 光滑小参数解的构造
% * 可测 Riemann 映照定理的准确引用
% * 复结构变化与坐标变化
% 
% 纲要内容结束

\section{Beltrami 方程与复结构}
\label{g4:sec:0602}

在固定光滑曲面上改变复结构，可以理解为改变哪些实切方向的复线性组合应当叫作 \((0,1)\) 方向。这样记录变化后，寻找新全纯坐标就成为一个一阶方程。本节先在光滑小参数范围内直接求解这个方程，并保留解对参数的依赖；下一节将把这些实际复结构中的坐标冗余除去。

\subsection{Beltrami 微分和局部变形}
\label{g4:subsec:0602:01}

固定一条光滑黎曼曲面 \(X\)，记 \(T_X=\mathrm{T}^{1,0}X=K_X^{-1}\) 为它的全纯切线丛。以下 \(A^{0,1}(T_X)\) 表示光滑 \(T_X\) 值 \((0,1)\) 形式。其元素在坐标 \(z\) 中写成
\[
	\mu=\mu_z\,\mathrm{d}\overline z\otimes\partial_z.
\]
若 \(\zeta=\zeta(z)\) 是原结构的另一个全纯坐标，则
\begin{equation}\label{g4:eq:0602:beltrami-transition}
	\mu_\zeta=\mu_z\frac{\zeta'(z)}{\overline{\zeta'(z)}}.
\end{equation}
所以 \(|\mu_z|\) 与坐标无关。满足 \(\|\mu\|_\infty<1\) 的这种张量称为\term[Beltramiweifen]{Beltrami 微分}[Beltrami differential]\footnote{Eugenio Beltrami（1835--1900），意大利数学家，研究微分几何，并用曲面模型解释非 Euclid 几何。}。在局部公式里，常把系数 \(\mu_z\) 也简写为 \(\mu\)。

\begin{proposition}\label{g4:prop:0602:almost-complex}
	设 \(X\) 是黎曼曲面，\(\mu\in A^{0,1}(T_X)\) 且 \(|\mu|<1\)。局部复直线
	\[
		\mathrm{T}^{0,1}_{\mu}X
		=\mathbb C(\partial_{\overline z}-\mu\partial_z)
	\]
	定义与原定向相容的光滑几乎复结构 \(J_\mu\)。相对于原结构足够近的每个几乎复结构都唯一这样表示。函数 \(w\) 是新结构的局部全纯坐标，当且仅当
	\begin{equation}\label{g4:eq:0602:beltrami-equation}
		w_{\overline z}=\mu w_z,\qquad w_z\ne0.
	\end{equation}
\end{proposition}
\begin{proof}
	式\eqref{g4:eq:0602:beltrami-transition}使所给直线在坐标交叠处相同。它与共轭直线的两个生成元的行列式为 \(1-|\mu|^2\)，所以二者互补。规定 \(J_\mu\) 在这两条直线上分别作用为 \(-\mathrm{i}\)、\(\mathrm{i}\)，即得到实切丛上的光滑端同态，满足 \(J_\mu^2=-I\)。其 \((1,0)\) 余切线由 \(\mathrm{d}z+\mu\,\mathrm{d}\overline z\) 张成；这个形式的实部、虚部确定的面积是原坐标面积的 \(1-|\mu|^2>0\) 倍，故定向相同。

	足够接近原 \((0,1)\) 线的复直线仍横截于原 \((1,0)\) 线，因而有唯一的图表示 \(\partial_{\overline z}-\mu\partial_z\)。最后，\(\mathrm{d}w\) 消灭新 \((0,1)\) 方向，恰为方程\eqref{g4:eq:0602:beltrami-equation}；实 Jacobian 为
	\[
		J_w=|w_z|^2-|w_{\overline z}|^2
		=(1-|\mu|^2)|w_z|^2.
	\]
	因此再加上 \(w_z\ne0\)，实逆函数定理就保证它确为保持定向的局部坐标。
\end{proof}

方程\eqref{g4:eq:0602:beltrami-equation}称为\term[Beltramifangcheng]{Beltrami 方程}[Beltrami equation]。写出几乎复结构还没有完成局部复图册的构造；必须找到满足方程且 Jacobian 非零的解。下一小节在本章使用的光滑小参数范围内完成这一步。

\begin{example}\label{g4:exm:0602:constant-beltrami}
	常数 \(c\in\mathbb C\)、\(|c|<1\) 对应显式坐标 \(w=z+c\overline z\)。若 \(c=re^{\mathrm{i}\theta}\)，适当旋转源、靶实坐标后，此实线性映射沿两条垂直方向的伸缩因子分别为 \(1+r\)、\(1-r\)。它把小圆变成椭圆，长短轴之比为 \((1+r)/(1-r)\)。这说明 \(|\mu|\) 衡量偏离共形的程度，而 \(\mu\) 的相位记录伸缩轴的方向。
\end{example}

\subsection{光滑小参数解的构造}
\label{g4:subsec:0602:02}

直接对方程作积分时，未知量又出现在右边的一阶导数中。要使迭代收敛，求解 \(\bar\partial\) 的算子必须恰好补回这一个导数。我们在一个辅助平环面上用 Fourier 系数构造这样的算子；环面的常数障碍由一个仿射项吸收，所求坐标只在原来的小圆盘上使用。

\begin{lemma}\label{g4:lem:0602:periodic-right-inverse}
	在 \(\mathbb T^2=\mathbb R^2/(2\pi\mathbb Z)^2\) 上记 \(H^s=W^{s,2}\)，\(M f=(4\pi^2)^{-1}\int_{\mathbb T^2}f\)。对每个整数 \(s\ge0\)，存在有界复线性算子
	\[
		T:H^s\longrightarrow H^{s+1},\qquad
		\bar\partial T f=f-Mf,\qquad MTf=0.
	\]
	这里 \(\bar\partial\) 在公式中指系数算子 \(\partial_{\overline z}\)。
\end{lemma}
\begin{proof}
	对 \(k=(k_1,k_2)\in\mathbb Z^2\) 取 Fourier 基 \(e^{\mathrm{i}(k_1x+k_2y)}\)，定义
	\[
		\widehat{Tf}(0)=0,\qquad
		\widehat{Tf}(k)=\frac{2\widehat f(k)}{\mathrm{i}k_1-k_2}
		\quad(k\ne0).
	\]
	分母的模为 \(|k|\)，所以
	\[
		\sum_k(1+|k|^2)^{s+1}|\widehat{Tf}(k)|^2
		\le8\sum_k(1+|k|^2)^s|\widehat f(k)|^2.
	\]
	这给出所需有界性；逐项乘以 \((\mathrm{i}k_1-k_2)/2\) 给出两个恒等式。光滑函数的级数可逐阶求导；一般情形由 Fourier 截断在 Sobolev 范数中的收敛得到。
\end{proof}

下面只固定一个整数 \(s\ge3\)。第三卷第11.3节推论11.3.10、11.3.11给出
\begin{equation}\label{g4:eq:0602:sobolev-tools}
	H^{s+1}\hookrightarrow C^1,
	\qquad \|fg\|_{H^s}\le C_s\|f\|_{H^s}\|g\|_{H^s}.
\end{equation}
这些结论在有限坐标覆盖和固定丛标架中逐分量适用。使用一个固定阶数，是为了先在完备的 Banach 空间\footnote{Stefan Banach（1892--1945），波兰数学家，建立完备赋范空间与线性算子的一般理论。}内求解；光滑性随后由方程提高，而不是在光滑函数的非 Banach 空间上直接使用逆函数定理。

\begin{theorem}\label{g4:thm:0602:small-beltrami}
	设 \(D\Subset U\subset\mathbb C\) 是圆盘，\(\mu(z,t)\) 在 \(U\) 上关于 \(z\) 光滑，参数 \(t\) 在 \(\mathbb C^N\) 中的原点附近变化，且 \(\mu(z,0)=0\)。假设 \(t\mapsto\mu(\cdot,t)\) 关于每个局部 Sobolev 范数全纯。缩小参数邻域后，\(D\) 上存在满足
	\[
		w_{\overline z}(z,t)=\mu(z,t)w_z(z,t),
		\qquad w(z,0)=z,
	\]
	的保持定向的光滑局部坐标。它关于 \(t\) 全纯，并在每个较小圆盘上关于空间变量和参数联合光滑。若参数依赖仅为实光滑，则所得坐标也实光滑地依赖参数。
\end{theorem}
\begin{proof}
	先经固定仿射缩放，使 \(\overline U\) 位于平环面的一个坐标方块内。取光滑截断 \(\chi\)，支于 \(U\) 且在 \(\overline D\) 附近等于 \(1\)，将 \(\chi\mu\) 在方块外补零并周期延拓，记为 \(m\)。这个过程保持参数的全纯性；\(m(0)=0\)，故 \(\|m(t)\|_{H^s}\) 在缩小参数域后任意小。

	寻找形如
	\[
		w(z,t)=z+a(t)\overline z+u(z,t),\qquad Mu=0
	\]
	的解，其中 \(u\) 周期。令
	\[
		A_m u=T(m\partial_z u),\qquad
		u=T m+A_m u.
	\]
	由式\eqref{g4:eq:0602:sobolev-tools}与上一引理，
	\[
		\|A_m u\|_{H^{s+1}}
		\le C\|m\|_{H^s}\|u\|_{H^{s+1}}.
	\]
	取 \(C\|m\|_{H^s}<1/2\)，则
	\begin{equation}\label{g4:eq:0602:neumann-solution}
		u=\sum_{j=0}^{\infty}A_m^jTm,
		\qquad a=M\bigl(m(1+\partial_z u)\bigr)
	\end{equation}
	在 \(H^{s+1}\) 中收敛，并且 \(\|u\|_{H^{s+1}}+|a|\le C'\|m\|_{H^s}\)。对方程施加 \(\bar\partial\)，得到
	\[
		\partial_{\overline z}u=m(1+\partial_z u)-a.
	\]
	于是 \(w_{\overline z}=m w_z\)，在 \(D\) 上即为所求方程。Sobolev 嵌入使 \(w\) 在 \(C^1\) 中接近 \(z\)，特别地 \(w_z\ne0\)。必要时再缩小参数域，\(\|\mathrm{d}w-I\|<1/2\) 于圆盘 \(D\)，沿线段积分还给出 \(|w(z)-w(z')|\ge|z-z'|/2\)，因此这里甚至可取单射坐标。

	每个级数项都由 \(m\)、复线性算子及双线性乘法构成，所以对 \(t\) 全纯。收敛在较小参数域上一致，故其和也为 Banach 值全纯函数；例如逐复参数应用 Cauchy 积分公式，并在一致收敛下交换积分与求和，就得到这一结论。实光滑参数的情形可对 \((I-A_m)^{-1}\) 求导，逆算子在这里一致有界，故得到各阶实参数导数。

	还需由有限阶解恢复光滑性。\(\partial_{\overline z}-m\partial_z\) 的实非零余切方向上的主符号不为零，因为两个复一阶符号的模相同而 \(|m|<1\)。它是光滑系数的一阶椭圆算子。第三卷第12.3节定理12.3.9及推论12.3.10的内部正则性用于 \(w\)，给出空间变量的光滑性。对参数微分方程，右端由 \(m\) 的参数导数和较低阶参数导数的 \(w_z\) 组成；在紧参数集上用相同的局部估计逐阶提高 Sobolev 阶数，再用嵌入，得到联合光滑性。参数域中的 Cauchy 公式也可在这些较高阶范数中使用，故全纯依赖与这项提高相容。
\end{proof}

\begin{corollary}\label{g4:cor:0602:holomorphic-family}
	设 \(X\) 为紧黎曼曲面，\(\eta_1,\ldots,\eta_N\in A^{0,1}(T_X)\)。对 \(t\in\mathbb C^N\) 足够小，令
	\[
		\mu(t)=\sum_{j=1}^N t_j\eta_j.
	\]
	则 \(J_{\mu(t)}\) 定义同一光滑曲面上的复结构 \(X_t\)。而且 \(X\times B\to B\) 可赋予复流形结构，使投影全纯、纤维为 \(X_t\)，其中 \(B\) 是足够小的参数球。
\end{corollary}
\begin{proof}
	取有限个内外嵌套坐标圆盘，在每个大盘上使用上一定理，在小盘上得到 \(w_i(z,t)\)。同一纤维上，两套坐标都消灭 \(\partial_{\overline z}-\mu(t)\partial_z\)，故其转移为全纯函数。还要核对参数方向：在 \((z,t)\) 的实光滑乘积上，新 \((0,1)\) 分布由
	\[
		\partial_{\overline z}-\mu(z,t)\partial_z,
		\quad \partial_{\overline t_1},\ldots,\partial_{\overline t_N}
	\]
	张成。函数 \(w_i(z,t),t_1,\ldots,t_N\) 全部消灭它，且其实 Jacobian 可逆。因此在坐标 \((w_i,t)\) 中，该分布就是标准 \((0,1)\) 分布；任意另一组坐标 \((w_j,t)\) 也消灭它，故转移满足全部 Cauchy--Riemann 方程而全纯。投影在这些坐标中就是 \((w_i,t)\mapsto t\)。
\end{proof}

这个构造已经产生实际复结构，不要求 Beltrami 微分满足附加的二阶可积条件。在复维一时，新反全纯方向只有一条复直线；这里用求解方程直接落实了它的局部坐标。

\subsection{可测 Riemann 映照定理的准确引用}
\label{g4:subsec:0602:03}

对一般可测系数，空间导数的逐阶估计不能使用，小参数的 Sobolev 压缩条件也未必成立。下面单独列出所需的标准外部存在定理。称保持定向的同胚 \(f\) 为 \(K\)-\term[zhungongxingyingshe]{准共形映射}[quasiconformal mapping]，是指它在局部坐标中属于 \(W^{1,2}_{\mathrm{loc}}\)，且几乎处处满足
\[
	|f_{\overline z}|\le\frac{K-1}{K+1}|f_z|,
	\qquad 1\le K<\infty.
\]
这与在几乎每条平行坐标轴的线段上绝对连续、并满足同一弱导数不等式的通常解析定义一致；绝对连续代表由一维 Sobolev 截面给出。

以下是\term[keceliemanyingzhaodingli]{可测 Riemann 映照定理}[measurable Riemann mapping theorem]的三点归一化形式。这里采用 Glutsyuk 论文作者稿的定理1.7及注1.8：定理1.7给出平面上固定 \(0,1\) 的解，注1.8保证无穷远处的准共形延拓。\cite[作者稿 arXiv:math/0510071，定理1.7、注1.8，第3页；证明见第3.1--3.4节]{Glutsyuk2008Uniformization}
\begin{theorem}[可测 Riemann 映照定理]\label{g4:thm:0602:measurable-riemann}
	设 \(\mu\in L^\infty(\mathbb C)\) 可测，\(\|\mu\|_\infty=k<1\)。则存在唯一保持定向的准共形同胚
	\[
		f:\widehat{\mathbb C\mkern0mu}\longrightarrow\widehat{\mathbb C\mkern0mu},
		\qquad f(0)=0,\quad f(1)=1,\quad f(\infty)=\infty,
	\]
	使 \(f_{\overline z}=\mu f_z\) 在 \(\mathbb C\) 上几乎处处成立；它的最大伸缩比至多为 \((1+k)/(1-k)\)。系数按几乎处处相等识别。这里不要求 \(\mu\) 连续或具有紧支集。若 \(\mu\) 在某个开集上光滑，则 \(f\) 在该开集上是光滑局部微分同胚。
\end{theorem}

三点归一化消除了后复合 Möbius 变换\footnote{August Ferdinand Möbius（1790--1868），德国数学家、天文学家，研究射影几何、仿射几何与拓扑。}的自由度。定理的结论包括全局单射和满射，不能从一个弱方程解的存在直接推出。本章的光滑局部参数由\cref{g4:thm:0602:small-beltrami}与下一节的切片构造证明；一般可测情形用于说明准共形标记所允许的范围。

下面同时记录规范化解的紧性，供比较准共形标记时使用。这是同一作者稿第3.1节的引理3.1、推论3.2及引理3.3；证明分别见第3.2、3.3节。\cite[作者稿第11--15页]{Glutsyuk2008Uniformization}

\begin{theorem}\label{g4:thm:0602:qc-compactness}
	固定 \(K\ge1\)。所有固定 \(0,1,\infty\) 的 \(K\)-准共形球面同胚及其逆映射，在球面距离下等度连续。因此任意这样的序列 \(f_n\) 都有子列一致收敛到一个球面同胚，逆映射同时一致收敛到其逆。

	更一般地，若平面域 \(U\) 上的 \(K\)-准共形同胚 \(f_n:U\to f_n(U)\) 在紧子集上一致收敛到同胚 \(f\)，且其 Beltrami 系数 \(\mu_n\) 几乎处处趋于 \(\mu\)，则 \(f\) 仍为 \(K\)-准共形并满足 \(f_{\overline z}=\mu f_z\)。特别地，若上述三点归一化球面映射的伸缩比 \(K_n\to1\)，则整个序列一致趋于恒等映射。
\end{theorem}

最后一句可直接从所引紧性与闭性推出：\(K_n\to1\) 使 \(\|\mu_n\|_\infty\to0\)，每个收敛子列的极限满足弱 Cauchy--Riemann 方程，所以全纯；三点归一化使它只能为恒等。若整个序列不趋于恒等，可取一个始终相距正数的子列，再取其收敛子列，就与这一唯一极限矛盾。一般可测系数的存在、全局单值性与上述紧性均是本节明确采用的外部输入；光滑小参数的构造和下一节的局部参数不依赖这些输入。

\subsection{复结构变化与坐标变化}
\label{g4:subsec:0602:04}

一个复结构可以用许多不同的标记坐标表示。比较两个 Beltrami 系数时，需要让坐标变换同时作用于源和靶，而不能只把两个系数相减。

\begin{proposition}\label{g4:prop:0602:change-of-structure}
	在固定光滑曲面上，设 \(J_\mu,J_\nu\) 由光滑 Beltrami 微分给出，\(\phi\) 是保持定向的光滑微分同胚。则 \(\phi:(X,J_\mu)\to(X,J_\nu)\) 全纯，当且仅当各原坐标中
	\begin{equation}\label{g4:eq:0602:map-equation}
		\phi_{\overline z}-\mu\phi_z
		+(\nu\circ\phi)
		\bigl(\overline{\phi_z}-\mu\overline{\phi_{\overline z}}\bigr)=0.
	\end{equation}
	若 \(w\) 是 \(J_{\mu_0}\) 的局部全纯坐标，则 \(J_\mu\) 相对于 \(w\) 的 Beltrami 系数为
	\begin{equation}\label{g4:eq:0602:relative-beltrami}
		\kappa
		=\frac{w_z}{\overline{w_z}}
		\frac{\mu-\mu_0}{1-\mu\overline{\mu_0}}.
	\end{equation}
\end{proposition}
\begin{proof}
	靶的新 \((1,0)\) 形式由 \(\mathrm{d}\zeta+\nu\,\mathrm{d}\overline\zeta\) 张成。它经 \(\phi\) 拉回后，\(\mathrm{d}z\)、\(\mathrm{d}\overline z\) 的系数分别为
	\[
		\phi_z+(\nu\circ\phi)\overline{\phi_{\overline z}},
		\qquad
		\phi_{\overline z}+(\nu\circ\phi)\overline{\phi_z}.
	\]
	全纯性要求后一系数等于前一系数乘 \(\mu\)，整理即得式\eqref{g4:eq:0602:map-equation}。作为符号检查，当 \(\phi=\operatorname{id}\) 时，该式恰给出 \(\mu=\nu\)。

	写 \(A=w_z\)，则 \(\mathrm{d}w=A(\mathrm{d}z+\mu_0\,\mathrm{d}\overline z)\)。解出 \(\mathrm{d}z,\mathrm{d}\overline z\) 并代入 \(\mathrm{d}z+\mu\,\mathrm{d}\overline z\)，得到
	\[
		\frac1{1-|\mu_0|^2}
		\left(
		\frac{1-\mu\overline{\mu_0}}A\,\mathrm{d}w
		+\frac{\mu-\mu_0}{\overline A}\,\mathrm{d}\overline w
		\right).
	\]
	两系数之比就是式\eqref{g4:eq:0602:relative-beltrami}；分母不为零，因为 \(|\mu\overline{\mu_0}|<1\)。
\end{proof}

式\eqref{g4:eq:0602:map-equation}特别适合下一节的局部切片：把源系数 \(\mu\) 视为给定数据，把靶系数 \(\nu\) 限制在一个有限维光滑空间，再求 \(\phi\)。这样只有光滑的有限维系数被复合到未知映射上，避免让任意有限阶数据 \(\mu\) 承担额外的空间导数。

\begin{example}\label{g4:exm:0602:torus-parameter}
	在 \(X_\tau=\mathbb C/(\mathbb Z+\tau\mathbb Z)\)、\(\operatorname{Im}\tau>0\) 上取常数系数 \(c\)、\(|c|<1\)。坐标 \(w=z+c\overline z\) 将两个周期变成 \(1+c\)、\(\tau+c\overline\tau\)。除以第一个周期后得到带相同有序标记的环面参数
	\[
		\tau(c)=\frac{\tau+c\overline\tau}{1+c},\qquad
		\operatorname{Im}\tau(c)
		=\frac{(1-|c|^2)\operatorname{Im}\tau}{|1+c|^2}>0.
	\]
	并且 \(\tau'(0)=\overline\tau-\tau\ne0\)。所以这个 Beltrami 参数确实改变了环面的复周期；局部坐标中的实线性伸缩，在整个商曲面上留下一个可计算的复模参数。
\end{example}
% 纲要条目排为编号小节标题，并进入目录与书签；说明段仍保留于纲要注释中。

% !TeX root = ../../Book4.tex
% 纲要标题：### 6.3 切空间、余切空间与维数
% 纲要位置：几何学计划纲要.md，第 3704 行。
% 正文状态：初稿；尚未正式审核。

% 纲要内容开始
% 
% * 平凡形变和向量场
% * Beltrami 类与上同调
% * 二次微分及 Serre 对偶
% * 稳定亏格情形的维数计算
% 
% 纲要内容结束

\section{切空间、余切空间与维数}
\label{g4:sec:0603}

上一节把一个小 Beltrami 微分变成了实际复结构，但不同微分可能只是同一曲面的不同坐标表达。我们先计算这种冗余的一阶形式，借助第三章的对偶找出不能由坐标变化消去的方向；随后求解一个非线性方程，证明这些方向确实形成局部参数。这样，维数计算所数的是附近曲面的参数，而不只是一个形式上的商空间。

\subsection{平凡形变和向量场}
\label{g4:subsec:0603:01}

仍固定紧连通黎曼曲面 \(X\)，记 \(T_X=K_X^{-1}\)。实向量场在全纯坐标中唯一写成
\[
	V=v\partial_z+\overline v\partial_{\overline z}.
\]
其 \((1,0)\) 分量 \(v\partial_z\) 是 \(T_X\) 的光滑截面；这给出实向量场与 \(A^0(T_X)\) 的实线性对应。令 \(\phi_t\) 是 \(V\) 的流，\(t\) 为实参数。紧性保证小时间内流在整个曲面上存在，且与恒等映射同位。

\begin{proposition}\label{g4:prop:0603:trivial-deformation}
	将固定复结构沿 \(\phi_t\) 拉回，并以原结构写其 Beltrami 系数 \(\mu(t)\)。则
	\[
		\mu(0)=0,\qquad
		\dot\mu(0)=\bar\partial_{T_X}(v\partial_z).
	\]
	因此由改变光滑标记产生的无穷小形变恰组成
	\(\bar\partial A^0(T_X)\subset A^{0,1}(T_X)\)。
\end{proposition}
\begin{proof}
	\(\phi_t:(X,\phi_t^*J)\to(X,J)\) 全纯。式\eqref{g4:eq:0602:map-equation}中取靶系数为零，得到
	\[
		\mu(t)=\frac{(\phi_t)_{\overline z}}{(\phi_t)_z}.
	\]
	局部 \(\phi_t(z)=z+t v(z)+O(t^2)\)，故分子为 \(t v_{\overline z}+O(t^2)\)，分母为 \(1+t v_z+O(t^2)\)，从而 \(\dot\mu(0)=v_{\overline z}\)。\(T_X\) 的标架 \(\partial_z\) 全纯，所以这正是 \(\bar\partial_{T_X}(v\partial_z)\) 的系数。任意光滑 \((1,0)\) 向量场均来自一个实向量场，故每个这样的像元素都由流实现。
\end{proof}

这里的“平凡”指一阶上由重标记产生，并非说系数函数为零。若一族复结构在一阶上满足 \(\dot\mu=\bar\partial v\)，只能立即断定它的一阶类为零；整族是否与常族等价，还需要非线性的坐标变化。

\begin{example}\label{g4:exm:0603:local-triviality}
	在一个坐标圆盘内取光滑紧支函数 \(v\)，将 \(v\partial_z\) 补零到整个曲面。非零的 \(v_{\overline z}\) 往往改变许多点上的复切方向，但它由上一命题的流产生，所以没有给出新的一阶模参数。另一方面，局部 \(\bar\partial\) 可解性说明每个 Beltrami 微分都能在一个小圆盘内这样消去；真正留下的信息在于这些局部消去能否粘成全局向量场。
\end{example}

\subsection{Beltrami 类与上同调}
\label{g4:subsec:0603:02}

定义无穷小形变空间
\begin{equation}\label{g4:eq:0603:infinitesimal-quotient}
	\mathcal V_X
	=\frac{A^{0,1}(T_X)}{\bar\partial A^0(T_X)}.
\end{equation}
分子记录所有一阶复结构变化，分母记录刚才算出的坐标变化。此时尚未由这个定义断言它是某个流形的切空间；这一点将在局部参数定理中得到证明。

\begin{proposition}\label{g4:prop:0603:dolbeault-class}
	对紧黎曼曲面 \(X\)，有自然同构
	\[
		\mathcal V_X\cong H^1(X,\mathcal O(T_X)).
	\]
	给定 \(\eta\in A^{0,1}(T_X)\)，在坐标覆盖上取 \(\bar\partial v_i=\eta\)。则它对应的全纯 Čech 余圈为 \(c_{ij}=v_i-v_j\)，其中差在同一切线丛内计算。其上同调类为零，当且仅当 \(\eta=\bar\partial v\) 有全局光滑解。
\end{proposition}
\begin{proof}
	\(T_X\) 是全纯线丛，故可直接将\cref{g4:thm:0302:dolbeault}用于 \(L=T_X\)。恢复该同构的局部意义：交叠处 \(\bar\partial(v_i-v_j)=0\)，所以 \(c_{ij}\) 是全纯向量场；三重交叠上差的相加等于零。改变局部解只增加全纯零上链的差，改变 \(\eta\) 为 \(\eta+\bar\partial v\) 则可同时把每个 \(v_i\) 改为 \(v_i+v\)，不改变余圈。

	若 \(\eta=\bar\partial v\)，则 \(v_i-v\) 全纯，故 \(c_{ij}\) 为余边界。反之，若 \(c_{ij}=a_i-a_j\)，其中 \(a_i\) 全纯，则 \(v_i-a_i\) 在交叠处相等，粘成所需的全局解。一般类的满射性与选择无关性已由第三章的单位分解构造证明。
\end{proof}

对于光滑参数族 \(\mu(t)\)，基点为 \(\mu(0)=0\)，把参数切向量 \(a\) 送到
\[
	\mathrm{KS}_0(a)=[\mathrm{d}\mu_0(a)]\in H^1(X,T_X)
\]
的映射称为该族在基点的\term[KodairaSpenceryingshe]{Kodaira--Spencer 映射}[Kodaira--Spencer map]\footnote{Kodaira，\OriginalHanName{小平邦彦}（\OriginalHanName{こだいら くにひこ}，Kunihiko Kodaira，1915--1997），日本数学家，研究复流形、代数曲面与形变理论。Spencer，Donald Clayton Spencer（1912--2001），美国数学家，研究复结构形变及微分方程的几何理论。}。改变族的光滑标记时，上一命题与\cref{g4:prop:0603:trivial-deformation}说明导数只增加 \(\bar\partial\) 向量场，故类不变。若纤维上的同构在基点不为恒等，则先用这个固定同构拉回比较，两边的类由它诱导的上同调同构对应。

\ProseParagraphBreak

在非零基点 \(\mu_0\) 处，须先把 \(X_{\mu_0}\) 当作新的基准曲面。由式\eqref{g4:eq:0602:relative-beltrami}，一阶变化 \(\eta\) 在它的全纯坐标 \(w\) 中变成
\begin{equation}\label{g4:eq:0603:relative-tangent}
	R_{\mu_0}\eta
	=\frac{w_z}{\overline{w_z}}
	\frac{\eta_z}{1-|\mu_0|^2}\,
	\mathrm{d}\overline w\otimes\partial_w.
\end{equation}
这是关于 \(\eta\) 的复线性变换。因而此处的 Kodaira--Spencer 类是 \([R_{\mu_0}\eta]\)。该公式也解释了为何不能在不同基点都把原来的系数差直接当作同一个上同调类。

\subsection{二次微分及 Serre 对偶}
\label{g4:subsec:0603:03}

与 Beltrami 类配对的对象应当抵消它的坐标变化因子。\(\eta\) 含有一个反全纯余切因子和一个全纯切因子；用两个全纯余切因子与它收缩后，恰剩一个可积分的 \((1,1)\) 形式。

\begin{definition}\label{g4:def:0603:quadratic-differential}
	黎曼曲面 \(X\) 上典范线丛的平方 \(K_X^{\otimes2}\) 的全纯截面称为\term[erciweifen]{全纯二次微分}[holomorphic quadratic differential]。其空间记为
	\[
		Q(X)=H^0(X,K_X^{\otimes2}).
	\]
	局部写作 \(q=q_z(\mathrm{d}z)^2\)，坐标变化满足 \(q_\zeta=q_z/(\zeta')^2\)。
\end{definition}

\begin{theorem}\label{g4:thm:0603:serre-duality}
	设 \(X\) 为紧黎曼曲面。积分
	\begin{equation}\label{g4:eq:0603:quadratic-pairing}
		\langle[\eta],q\rangle
		=\frac{\mathrm{i}}2\int_X q\mathbin{\lrcorner}\eta
		=\int_Xq_z\eta_z\,\mathrm{d}x\,\mathrm{d}y
	\end{equation}
	给出 \(\mathcal V_X\times Q(X)\) 上非退化的复双线性配对。第一处积分中，收缩后的形式按 \(\mathrm{d}z\wedge\mathrm{d}\overline z\) 排列；末式按坐标分片相加理解。因此有自然复线性同构
	\[
		\mathcal V_X^\vee\cong Q(X).
	\]
	若固定共形度量 \(g_0=\rho^2|\mathrm{d}z|^2\)，并在 \(T_X\) 上取局部标架 \(\partial_z\) 的平方范数为 \(\rho^2\)，则每个 Beltrami 类具有唯一形如
	\begin{equation}\label{g4:eq:0603:harmonic-beltrami}
		\eta=\rho^{-2}\overline{q_z}\,
		\mathrm{d}\overline z\otimes\partial_z,
		\qquad q\in Q(X),
	\end{equation}
	的调和代表。
\end{theorem}
\begin{proof}
	坐标变换时，\(q_z\eta_z\) 乘以 \(|\zeta'|^{-2}\)，而面积元乘以 \(|\zeta'|^2\)，所以被积二形式全局有定义。将\cref{g4:thm:0303:serre-duality}用于 \(L=T_X=K_X^{-1}\)，其对偶线丛为
	\[
		K_X\otimes T_X^{-1}=K_X^{\otimes2}.
	\]
	第三章的配对常数为 \(1/(2\pi\mathrm{i})\)，式\eqref{g4:eq:0603:quadratic-pairing}等于它乘以 \(-\pi\)。非零常数不改变非退化性与自然同构。也可直接看出代表无关：若 \(\eta=\bar\partial v\)，则 \(q\mathbin{\lrcorner}\bar\partial v\) 是全局一形式外微分的负值，紧无边界曲面上的 Stokes 公式使其积分为零。

	再恢复调和代表的公式。\cref{g4:prop:0303:formal-adjoint}在所取切线丛度量下给出
	\[
		\bar\partial_{T_X}^*\eta=0
		\quad\Longleftrightarrow\quad
		\partial_z(\rho^2\eta_z)=0.
	\]
	所以 \(\rho^2\overline{\eta_z}\) 是全纯函数，并且按二次微分的系数变换，得到式\eqref{g4:eq:0603:harmonic-beltrami}。反过来，该式的每个截面均满足伴随方程。\cref{g4:thm:0303:fredholm}的正交分解保证每个类中这样的代表存在且唯一。若 \(\eta\ne0\) 调和，取 \(q_z=\rho^2\overline{\eta_z}\)，则
	\[
		\langle[\eta],q\rangle
		=\int_X\rho^2|\eta_z|^2\,\mathrm{d}x\,\mathrm{d}y>0.
	\]
	这也直接显示了配对如何检测非零类。式\eqref{g4:eq:0603:harmonic-beltrami}的对应为反线性，配对给出的余切识别则为复线性，二者的类型应予区别。
\end{proof}

记上述有限维调和空间为 \(\mathscr H_X\)。于是
\begin{equation}\label{g4:eq:0603:hodge-splitting}
	A^{0,1}(T_X)=\bar\partial A^0(T_X)\oplus\mathscr H_X.
\end{equation}
选择度量只选择了每个类的一个方便代表；\(H^1(X,T_X)\)、\(Q(X)\) 及积分配对本身不依赖这个选择。下一小节将把 \(\mathscr H_X\) 作为实际复结构族的参数方向。

\subsection{稳定亏格情形的维数计算}
\label{g4:subsec:0603:04}

现在假设 \(g=g(X)\ge2\)。这一条件同时带来两项简化：全纯向量场为零，故近恒等的重标记没有无穷小稳定子；二次微分的次数足够高，Riemann--Roch 中的对偶项消失。

\begin{corollary}\label{g4:cor:0603:dimension}
	设 \(X\) 是亏格 \(g\ge2\) 的紧连通黎曼曲面。则
	\[
		H^0(X,T_X)=0,\qquad
		\dim_{\mathbb C}Q(X)
		=\dim_{\mathbb C}\mathcal V_X=3g-3.
	\]
\end{corollary}
\begin{proof}
	由\cref{g4:cor:0304:canonical-degree,g4:thm:0304:genus-comparison}，
	\[
		\deg T_X=2-2g<0,
		\qquad \deg K_X^{\otimes2}=4g-4.
	\]
	负次数线丛没有非零全纯截面，因为非零全纯截面的零除子有效，次数不能为负。这证明第一项。将 Riemann--Roch 用于 \(L=K_X^{\otimes2}\)，其对偶项是 \(H^0(X,K_X^{-1})=0\)，所以
	\[
		h^0(K_X^{\otimes2})=4g-4+1-g=3g-3.
	\]
	再用\cref{g4:thm:0603:serre-duality}得到 \(\mathcal V_X\) 的维数。
\end{proof}

下面证明这些方向真能作局部参数。为让无限维求解中每一步都有确定的空间，固定整数 \(s\ge3\)，所有 Sobolev 范数使用 \(X\) 上固定的有限图册、丛标架和光滑度量。

\begin{lemma}\label{g4:lem:0603:split-inverse}
	在上述 \(g\ge2\) 条件下，算子
	\[
		\mathcal L:H^{s+1}(T_X)\oplus\mathscr H_X
		\longrightarrow H^s(\Lambda^{0,1}\otimes T_X),
		\qquad \mathcal L(v,h)=\bar\partial v+h
	\]
	是有界复线性同构。
\end{lemma}
\begin{proof}
	\(\bar\partial_{T_X}\) 是两个复线丛间的一阶椭圆算子。第三章的 \(L^2\) 分解给出右端 \(b\) 的唯一调和分量 \(h\)；剩余部分与伴随核正交。第三卷第13.2节推论13.2.5将其解提高到 \(H^{s+1}\)，并给出
	\[
		\|v\|_{H^{s+1}}\le C_s\|b-h\|_{H^s}.
	\]
	这里无需再规定核正交归一化，因为 \(\ker\bar\partial=H^0(X,T_X)=0\)。调和投影是由有限个光滑基向量给出的有界有限秩算子，所以 \(h\) 的范数也受 \(\|b\|_{H^s}\) 控制。存在性、唯一性和有界逆因此全部成立。
\end{proof}

我们将用一次 Banach 空间隐函数定理，其求解机制可简述如下。若光滑映射 \(F(x,p)\) 满足 \(F(0,0)=0\)，且 \(A=D_xF(0,0)\) 有有界逆，则
\[
	x\longmapsto x-A^{-1}F(x,p)
\]
在一个足够小的闭球上具有小于 \(1/2\) 的 Lipschitz 常数\footnote{Rudolf Otto Sigismund Lipschitz（1832--1903），德国数学家，研究微分方程、微分几何与数论。}；缩小 \(p\) 的邻域使它还将此球送入自身。逐次迭代为 Cauchy 列，给出唯一零点。两参数的方程相减先得解连续；再用差商和 \(D_xF\) 的可逆性得到
\(D_px=-(D_xF)^{-1}D_pF\)，逐阶微分得到光滑依赖。以下会明确写出 \(F\)、空间和可逆的导数，而不在光滑微分同胚群上直接使用这一结论。

\begin{theorem}\label{g4:thm:0603:local-parameters}
	设 \(S\) 为闭连通定向光滑曲面，亏格 \(g\ge2\)，\((X,f)\) 为其一个光滑标记的黎曼曲面。取 \(\mathscr H_X\) 的复基 \(\eta_1,\ldots,\eta_{3g-3}\)，令
	\[
		\nu(t)=\sum_{j=1}^{3g-3}t_j\eta_j.
	\]
	存在 \(0\in B\subset\mathbb C^{3g-3}\) 的开邻域，使
	\begin{equation}\label{g4:eq:0603:teich-chart}
		B\longrightarrow\mathcal T(S),\qquad
		t\longmapsto[X_{\nu(t)},f]
	\end{equation}
	是到 \([X,f]\) 的一个开邻域的同胚，其中 \(\mathcal T(S)\) 采用光滑复结构模同位于恒等的微分同胚之商拓扑。不同基准曲面及调和基给出的这些图的转移映射全纯。因此 \(\mathcal T(S)\) 是局部复维数 \(3g-3\) 的复流形，并且
	\[
		\mathrm T_{[X,f]}\mathcal T(S)\cong H^1(X,T_X),
		\qquad
		\mathrm T^*_{[X,f]}\mathcal T(S)\cong Q(X).
	\]
\end{theorem}
\begin{proof}
	先求一个近恒等的规范化。把 \(X\) 的底层光滑曲面固定，给它取辅助光滑黎曼度量。对小的 \(v\in H^{s+1}(T_X)\)，将其与共轭相加看成实向量场，令 \(\phi_v(x)=\exp_x(V(x))\)。由 \(H^{s+1}\hookrightarrow C^1\)，\(v\) 小时 \(\phi_v\) 是 \(H^{s+1}\) 类的 \(C^1\) 微分同胚，且 \(r\mapsto\phi_{rv}\) 给出到恒等的 \(C^1\) 同位。为证明全局单射，可先由 \(C^1\) 接近恒等得到局部微分同胚；它在紧曲面上是覆盖，而这条同伦使其度数为 \(1\)，故为一层覆盖。

	给定小的源系数 \(\mu\in H^s(\Lambda^{0,1}\otimes T_X)\)，寻找 \(v\) 和 \(h\in\mathscr H_X\)，使 \(\phi_v:(X,J_\mu)\to(X,J_h)\) 全纯。局部残差正是
	\begin{equation}\label{g4:eq:0603:slice-equation}
		E(v,h,\mu)
		=(\phi_v)_{\overline z}-\mu(\phi_v)_z
		+(h\circ\phi_v)
		\bigl(\overline{(\phi_v)_z}
		-\mu\overline{(\phi_v)_{\overline z}}\bigr).
	\end{equation}
	它本来取值于沿 \(\phi_v\) 的靶切线丛；用固定联络沿短的指数曲线平行移动，将它识别为固定丛 \(\Lambda^{0,1}\otimes T_X\) 的截面。这个识别在 \(v=0\) 为恒等，且光滑依赖 \(x,v(x)\)，故不改变下面的零点条件和线性化。

	核对函数空间是这一构造的关键。未知 \(v\) 在 \(H^{s+1}\)，它的一阶导数在 \(H^s\)；\(\mu\) 仅与这些导数相乘。被复合的 \(h\) 则始终在有限维的光滑空间 \(\mathscr H_X\) 内。固定光滑函数的代入 \(v\mapsto a(x,v(x))\) 在这里光滑：对 \(a\) 用带积分余项的 Taylor 公式，再逐阶求空间导数，各项由式\eqref{g4:eq:0602:sobolev-tools}的乘法估计控制；其各阶微分就是相应的偏导数乘以增量。有限个坐标片求和给出相同结论。因此
	\[
		E:H^{s+1}(T_X)\times\mathscr H_X\times H^s(\Lambda^{0,1}\otimes T_X)
		\longrightarrow H^s(\Lambda^{0,1}\otimes T_X)
	\]
	在原点附近是实光滑映射。这里没有出现 \(\mu\circ\phi_v\)：如果直接对任意 \(H^s\) 系数作这种复合再对 \(v\) 求导，便会出现未受控的 \(\mathrm{d}\mu\)，不能得到同一阶空间中的上述结论。

	在原点微分，得到
	\[
		D E_{(0,0,0)}(v,h,\eta)=\bar\partial v+h-\eta.
	\]
	\cref{g4:lem:0603:split-inverse}说明关于 \((v,h)\) 的导数可逆。刚才说明的隐函数论证于是给出唯一的小解 \((v(\mu),P(\mu))\)，实光滑地依赖 \(\mu\)。若 \(\mu\) 光滑，则源、靶由\cref{g4:thm:0602:small-beltrami}具有光滑局部复坐标；\(\phi_v\) 在这些坐标中是满足 Cauchy--Riemann 方程的 \(C^1\) 映射，故全纯并光滑。于是每个足够接近 \(X\) 的光滑复结构，均经一个同位于恒等的光滑微分同胚，成为唯一的小调和系数候选。

	接着把“近恒等规范化唯一”加强为标记等价类中的唯一。对此需要排除两个趋近原点的调和系数之间存在离恒等很远的同位映射。先说明有限维族 \(X_h\) 的双曲度量随 \(h\) 光滑变化。设 \(g_0\) 为 \(X\) 的 Poincaré 度量，取
	\[
		b_h(U,W)=\tfrac12\bigl(g_0(U,W)+g_0(J_hU,J_hW)\bigr).
	\]
	它是 \(J_h\) 的光滑共形度量，\(b_0=g_0\)。由式\eqref{g4:eq:0504:curvature-change}，求曲率 \(-1\) 的 \(e^{2u}b_h\) 等价于
	\[
		F(u,h)=\Delta_{b_h}u+K_{b_h}+e^{2u}=0.
	\]
	这里 \(F:H^{s+2}(X,\mathbb R)\times\mathscr H_X\to H^s(X,\mathbb R)\) 实光滑，关于 \(u\) 的原点导数为 \(\Delta_{g_0}+2\)。它的核为零，因为与 \(u\) 配对得到 \(\|\mathrm{d}u\|_2^2+2\|u\|_2^2\)；第三卷第13.2节定理13.2.4及推论13.2.5给出满射和有界逆。隐函数论证得到 \(u(h)\)。方程的椭圆估计逐阶提高正则性，对参数求导后也同样成立，故 \(g_h=e^{2u(h)}b_h\) 关于空间与 \(h\) 光滑，且 \(g_h\to g_0\) 于每个固定 \(C^k\) 范数。

	同一紧共形类的曲率 \(-1\) 度量唯一。事实上，若另一个为 \(e^{2a}g_h\)，则 \(\Delta_{g_h}a=1-e^{2a}\)。正的最大值处，非负 Laplace 的左边非负，右边负，矛盾；负的最小值同样不可能。因此 \(a=0\)。特别地，任何双全纯映射 \(\psi:X_h\to X_{h'}\) 都是 \(g_h,g_{h'}\) 之间的等距映射。

	若任意小邻域中仍有不同 \(h_n,h'_n\to0\) 由同位于恒等的双全纯映射 \(\psi_n\) 联系，等距等式和 \(g_{h_n},g_{h'_n}\to g_0\) 给出 \(\mathrm{d}\psi_n\) 与逆映射微分的一致界。等距映射将靶的 Levi-Civita 联络拉回为源的度量相容无挠联络，由该联络的唯一性知它保持联络，因此坐标中
	\[
		\partial_i\partial_j\psi_n^a
		=\Gamma^k_{ij}(g_{h_n})\partial_k\psi_n^a
		-\Gamma^a_{bc}(g_{h'_n})(\psi_n)
		\partial_i\psi_n^b\partial_j\psi_n^c.
	\]
	有限图册内先由此控制二阶导数，再反复求导控制全部高阶导数。对映射及逆映射用紧性和对角选子列，得到 \(C^\infty\) 收敛到 \(X\) 的全纯自同构 \(\psi\)。足够一致接近的两个映射可沿短测地线连接成同伦，所以 \(\psi\) 同伦于恒等。\cref{g4:prop:0601:automorphism-faithful}给出 \(\psi=\operatorname{id}\)。因而最终 \(\psi_n\) 属于刚才隐函数论证的近恒等邻域。对于源系数 \(h_n\)，一组解是 \((v,h)=(0,h_n)\)，另一组是 \((\exp^{-1}\psi_n,h'_n)\)，唯一性强迫 \(h'_n=h_n\)，矛盾。这就得到所需的局部单射性。

	实光滑的映射 \(P\) 在光滑系数的 \(C^\infty\) 拓扑中连续，而 \(P(h)=h\) 对小调和系数成立。它在刚才的小邻域内恰按标记等价类取相同值。因此 \(h\mapsto[X_h,f]\) 与 \(P\) 互为连续的局部逆。这个像是开邻域：一个小的系数开集在全部同位重标记下的饱和仍开，因为每个固定光滑微分同胚的拉回都是 \(C^\infty\) 拓扑的同胚；在任意饱和点附近，先用该固定拉回移回小系数邻域，再用 \(P\)。缩小调和参数球即可得到式\eqref{g4:eq:0603:teich-chart}的拓扑结论。

	最后证明参数转移全纯，这一步不能由实隐函数定理自动得到。由\cref{g4:cor:0602:holomorphic-family}，\(h(t)=\sum t_j\eta_j\) 给出的族本身是全纯曲面族。在任意小参数 \(t_0\) 处，其 Kodaira--Spencer 映射为
	\[
		\mathbb C^{3g-3}\longrightarrow H^1(X_{h(t_0)},T_{X_{h(t_0)}}),
		\qquad a\longmapsto[R_{h(t_0)}h(a)].
	\]
	它复线性，并且在缩小参数域后处处可逆。具体地，在 \((v,h,\mu)=(0,h_0,h_0)\) 处，\(D_{(v,h)}E\) 是原点同构的算子范数小扰动，故仍可逆；其逆由 \((I+A^{-1}\delta A)^{-1}A^{-1}\) 的 算子几何级数给出。改用 \(X_{h_0}\) 的全纯坐标和相应的 \((1,0)\) 向量场变量后，该导数的零点与
	\[
		\bar\partial_{X_{h_0}}V+R_{h_0}h(a)=0
	\]
	的零点相同。可逆性说明这些参数方向与 \(\bar\partial\) 的像互补；取商后恰得到上述 Kodaira--Spencer 同构。

	设另一组这样的参数记为 \(s\)，重叠图的转移为 \(s=s(t)\)。固定重叠点处的标记同构后，隐函数构造给出局部实光滑的比较映射及 \(s(t)\)。沿任意 \(\partial/\partial\overline t_j\) 微分两族的标记等价关系。源族的 Beltrami 系数关于 \(t\) 全纯，因此在此方向的 Kodaira--Spencer 类为零。靶族的类为
	\[
		\mathrm{KS}_{s(t)}\left(\frac{\partial s}{\partial\overline t_j}\right).
	\]
	比较映射的变化只增加 \(\bar\partial\) 向量场；这是\cref{g4:prop:0603:trivial-deformation}在当前纤维的坐标中求导后，再对实参数方向复线性延拓的结论。故上式也为零。刚证的 Kodaira--Spencer 可逆性给出 \(\partial s/\partial\overline t_j=0\)。实光滑转移满足 Cauchy--Riemann 方程，因而全纯；反向转移同理。

	这些局部图在每个标记曲面处都可构造。局部 Hausdorff 性由参数球得到；不同标记等价类若不能分离，取两个基准曲面的系数邻域中相互等价的序列，同样的双曲等距紧性给出它们之间的极限全纯同构，而在 \(C^1\) 意义下足够接近的微分同胚，可沿短测地线插值成一条仍为微分同胚的路径，所以极限保留标记的同位类，产生矛盾。因此商空间 Hausdorff。光滑张量空间有可数拓扑基，商映射开，所以商也第二可数。这些图遂构成复流形图册。其基点导数正是调和代表到 \(H^1(X,T_X)\) 的同构，余切识别再由\cref{g4:thm:0603:serre-duality}得到。
\end{proof}

这一定理的证明分开了三个问题：椭圆分解给出可逆线性化，非线性映射方程产生实际参数，参数族的全纯性保证复图册。

\begin{example}\label{g4:exm:0603:low-genus}
	亏格 \(0\) 时 \(T_{\mathbb P^1}\) 次数为 \(2\)，全纯向量场三维，\(Q(\mathbb P^1)=0\)。亏格 \(1\) 时 \(T_X\) 平凡，两个空间 \(H^0(X,T_X)\)、\(Q(X)\) 均为一维；\cref{g4:exm:0602:torus-parameter}已给出其实际一维复参数。低亏格中的全纯向量场使\cref{g4:lem:0603:split-inverse}的核不再为零，须先固定适当的点或平移归一化。对 \(g\ge2\)，这个无穷小核消失，但仍可能存在离散的全纯自同构；它们改变标记，在忘掉标记的模空间中还要另取有限商。
\end{example}
% 纲要条目排为编号小节标题，并进入目录与书签；说明段仍保留于纲要注释中。

% !TeX root = ../../Book4.tex
% 纲要标题：### 6.4 实际参数与全局几何
% 纲要位置：几何学计划纲要.md，第 3711 行。
% 正文状态：初稿；尚未正式审核。
% 本次扩充的范围、先修与边界：先修为 IV.5 的双曲统一化与测地几何、IV.6.1 的标记和 Dehn 扭转。主对象为亏格至少二的闭连通定向双曲曲面；选定带标记的裤子分解和缝线，说明长度、扭转的单位、参数数量、有限拼接和恢复。长度趋零只处理选定曲线捏缩的模型，区别 Teichmüller 空间、模空间与稳定退化；不追加完整 Weil–Petersson 几何、Wolpert 公式或一般全局紧化理论。

% 纲要内容开始
% 
% > 补充写作范围与外部输入：先修为 IV.5 的双曲统一化与测地几何、IV.6.1 的标记和 Dehn 扭转。主对象为亏格至少二的闭连通定向双曲曲面；选定带标记的裤子分解和缝线，说明长度、扭转的单位、参数数量、有限拼接和恢复。长度趋零只处理选定曲线捏缩的模型，区别 Teichmüller 空间、模空间与稳定退化；不追加完整 Weil–Petersson 几何、Wolpert 公式或一般全局紧化理论。
% 
% * 局部坐标及自同构的影响
% * 极值长度和 Teichmüller 度量
% * Fenchel–Nielsen 参数的模型
% * 退化、节点和紧化的入口
% 
% * 双曲裤子的边界长度与裤子分解
% * 长度、扭转与标记的选择
% * 沿测地边界的拼接与参数恢复
% * 长度趋零时的标记与退化
% 
% 
% 纲要内容结束

\section{实际参数与全局几何}
\label{g4:sec:0604}

切空间的维数回答了小形变有多少个独立方向；实际参数还须告诉我们如何制造曲面，以及怎样从曲面读回原来的数值。局部复坐标由上一节的切片给出，本节再用双曲裤子的有限拼接构造全局实参数。与此相配，极值长度比较度量形变的大小，长度趋零的模型则说明这些参数怎样走向光滑曲面的边界。

\subsection{局部坐标及自同构的影响}
\label{g4:subsec:0604:01}

固定闭连通定向曲面 \(S\)，亏格 \(g\ge2\)。\cref{g4:thm:0603:local-parameters}把标记黎曼曲面 \([X,f]\) 的一个邻域表示为 \(\mathbb C^{3g-3}\) 中的开集：有限维参数来自 Beltrami 类，光滑坐标变化的方向由 \(\bar\partial\) 像消去。这里的参数点仍带标记；忘记标记以后，还须识别由曲面自同构及映射类作用产生的等价。

\begin{proposition}
	\label{g4:prop:0604:finite-automorphisms}
	设 \(X\) 为亏格 \(g\ge2\) 的紧连通黎曼曲面，则 \(\operatorname{Aut}(X)\) 有限。
\end{proposition}
\begin{proof}
	每个自同构保持 \(X\) 的 Poincaré 度量，这是\cref{g4:thm:0504:poincare-complete}对映射及其逆同时应用的结果。于是自同构族在紧度量空间上等度连续；取有限网格并作对角选择，任意序列都存在一致收敛子列。对逆映射作同样选择，极限互为逆，因而是等距同胚。局部取相对紧全纯坐标片，一致收敛与 Cauchy 公式又说明极限及其逆全纯。因此自同构族在一致拓扑下紧。

	若它无限，可以由上述紧性构造互不相同的 \(F_n\to\operatorname{id}\)。记 \(\varepsilon_n=\max_{p\in X}d_X(F_n(p),p)>0\)。在有限张较小坐标片内，\((F_n(z)-z)/\varepsilon_n\) 局部一致有界；由 Montel 定理同时取子列，使它们局部一致收敛。若换坐标 \(w=h(z)\)，Taylor 展开给出
	\[
		\frac{h(F_n(z))-h(z)}{\varepsilon_n}
		=h'(z)\frac{F_n(z)-z}{\varepsilon_n}+O(\varepsilon_n).
	\]
	故极限按向量场的变换律粘合为一个全纯向量场 \(V\)。选取实现最大位移的点并再取收敛子列，局部度量与坐标距离的双边比较说明 \(V\ne0\)。但 \(\deg T_X=2-2g<0\)，由\cref{g4:thm:0305:high-degree-sections}，负次线丛没有非零全纯截面，矛盾。因此自同构群有限。
\end{proof}

局部模空间确实只需作这个有限商。为说明不会有其他映射类从很远处回到同一小邻域，取局部参数 \(a\) 对应的光滑双曲度量 \(h_a\)，它在固定标记曲面上随 \(a\) 光滑变化。假如有 \(a_n,b_n\to0\) 及等价 \(F_n:(S,h_{a_n})\to(S,h_{b_n})\)，则这些等距映射及其逆在固定背景度量下具有统一的 Lipschitz 界。上面的紧性论证给出子列收敛于 \((S,h_0)\) 的自同构 \(F\)。还可把收敛提高到光滑收敛：在一点选取微分的收敛子列，局部等距满足 \(F_n(\exp_p^{h_{a_n}}v)=\exp_{F_n(p)}^{h_{b_n}}(\mathrm{d}F_n|_p v)\)，而测地线方程随初值及光滑度量光滑变化。有限个小片逐次延续便得到全曲面上的 \(C^1\) 收敛。当 \(F^{-1}F_n\) 在 \(C^1\) 意义下足够接近恒等时，以短测地线写成 \(p\mapsto\exp_pV_n(p)\)，并用 \(p\mapsto\exp_p(tV_n(p))\)、\(0\le t\le1\) 插值。每一步的微分都可逆，紧性使它为有限覆盖，同伦度为 \(1\) 又使它为微分同胚，故这确实是一条同位。因此 \(F_n\) 最终属于 \(F\) 的映射类。

\ProseParagraphBreak

由于 \(\operatorname{Aut}(X)\) 有限，可缩小局部参数邻域，并对有限群取共同不变的邻域，使忘记标记的局部空间恰是该邻域模有限稳定子。作用是全纯的，因为标记变换与局部切片的坐标转移均为全纯。有限群作用还可局部线性化：若各 \(A\) 固定原点，置
\[
	\Psi(z)=\frac1{|G|}\sum_{A\in G}(\mathrm{d}A_0)^{-1}A(z).
\]
则 \(\mathrm{d}\Psi_0=I\)，且 \(\Psi(Az)=\mathrm{d}A_0\Psi(z)\)，逆函数定理给出所需坐标。因此局部模空间可写成一个有限线性群商。该表示可能有核；非平凡稳定子也不自动造成商空间的奇点。

\subsection{极值长度和 Teichmüller 度量}
\label{g4:subsec:0604:02}

设 \(F:X\to Y\) 为保向 \(C^1\) 微分同胚。在任一点取共形正交坐标，令其微分的两个奇异值为 \(s_1\ge s_2>0\)，定义最大伸缩比
\[
	K(F)=\sup_X\frac{s_1}{s_2}\ge1.
\]
这个量不受源、目标共形密度的改变影响；在复坐标中它等于 \((1+\|\mu_F\|_\infty)/(1-\|\mu_F\|_\infty)\)。当 \(K(F)=1\) 时，微分是复线性的，故 \(F\) 双全纯。

\begin{proposition}
	\label{g4:prop:0604:extremal-qc}
	设 \(F:X\to Y\) 是保向 \(C^1\) 微分同胚，且 \(K(F)\le K<\infty\)。对 \(X\) 上任一曲线族 \(\Gamma\)，有
	\[
		K^{-1}\operatorname{EL}_X(\Gamma)
		\le\operatorname{EL}_Y(F\Gamma)
		\le K\operatorname{EL}_X(\Gamma).
	\]
\end{proposition}
\begin{proof}
	给定目标可容许密度 \(\sigma\)，在源上取 \(\rho=(\sigma\circ F)s_1\)。链式法则说明每条曲线都有
	\[
		L_\sigma(F\gamma)\le L_\rho(\gamma).
	\]
	另一方面，\(s_1^2\le K s_1s_2=KJ_F\)，面积换元给出
	\[
		A_\rho\le K\int_X(\sigma\circ F)^2J_F\,\mathrm{d}A_X
		=KA_\sigma.
	\]
	因此 \(L_\sigma(F\Gamma)^2/A_\sigma\le K\operatorname{EL}_X(\Gamma)\)。对 \(\sigma\) 取上确界得到右边不等式；对 \(F^{-1}\) 应用同样论证得到左边。背景共形度量的选择不影响计算。
\end{proof}

一般准共形映射只有弱导数，沿某条指定曲线未必绝对连续；因而不能逐条直接套用上述证明。本节对这种映射使用环域版本，证明中只沿几乎每条平行圆周积分。这里所需的逆映射、复合、零测集保持及面积换元性质，采用准共形解析理论的标准结论；下面的证明也是 Glutsyuk 作者稿引理3.7的圆柱形式。\cite[作者稿第12页，引理3.7；群性质与零测集保持见第3.4节]{Glutsyuk2008Uniformization}

\begin{proposition}
	\label{g4:prop:0604:annulus-qc}
	设 \(A_1,A_2\) 为有限模量的共形圆环域，\(F:A_1\to A_2\) 为 \(K\)-准共形同胚，则
	\[
		K^{-1}\operatorname{Mod}(A_1)
		\le\operatorname{Mod}(A_2)
		\le K\operatorname{Mod}(A_1).
	\]
	因此环域的连接族与核心族也分别满足同一极值长度比较。
\end{proposition}
\begin{proof}
	利用共形坐标，将 \(A_j\) 写成平坦圆柱
	\[
		(0,L_j)\times(\mathbb R/2\pi\mathbb Z),
		\qquad \operatorname{Mod}(A_j)=L_j/(2\pi).
	\]
	记源坐标为 \((x,\theta)\)。局部 \(W^{1,2}\) 性说明对几乎所有 \(x\)，\(F(x,\theta)\) 关于 \(\theta\) 绝对连续。这个闭路在目标圆柱的绕行数为 \(\pm1\)，所以其平坦长度至少为 \(2\pi\)。几乎处处的伸缩比界给出
	\[
		|\partial_\theta F|^2\le KJ_F.
	\]
	先在 \(\varepsilon<x<L_1-\varepsilon\) 积分，再令 \(\varepsilon\downarrow0\)，由面积换元、Fubini 定理\footnote{Guido Fubini（1879--1943），意大利数学家，研究分析、微分几何与重积分。}与 Cauchy--Schwarz 不等式，
	\[
		2\pi L_2=\int_{A_1}J_F
		\ge\frac1K\int_0^{L_1}\int_0^{2\pi}|\partial_\theta F|^2
		\,\mathrm{d}\theta\,\mathrm{d}x
		\ge\frac{2\pi L_1}{K}.
	\]
	这给出左边不等式。逆映射仍为 \(K\)-准共形，交换源、目标便得右边。最后由\cref{g4:thm:0505:annulus-modulus}，两族的极值长度分别等于模量及其倒数，所以得到所述比较。
\end{proof}

标记规定了允许比较的映射同伦类。\term[Teichmullerjuli]{Teichmüller 距离}[Teichmüller distance]定义为
\begin{equation}
	\label{g4:eq:0604:teich-distance}
	d_{\mathrm T}([X,f],[Y,h])
	=\frac12\inf_{F\simeq h\circ f^{-1}}\log K(F),
\end{equation}
其中下确界取遍该类的准共形同胚；同伦类与光滑同位标记的衔接由\cref{g4:prop:0601:surface-isotopy-interface}保证。\(K(F)\) 是几乎处处伸缩比的本质上确界。取一个光滑标记比较映射，紧性便给出有限的 \(K(F)\)，所以距离值有限。逆映射保持 \(K\)，复合满足 \(K(GF)\le K(G)K(F)\)，因而对称性与三角不等式直接来自伸缩比的运算。

\ProseParagraphBreak

严格正定性及完备性涉及准共形映射族的全局紧性。这里把以下结论列为外部输入：式\eqref{g4:eq:0604:teich-distance}确实是 \(\mathcal T(S)\) 上的完备距离，并由其通常复流形切空间上的连续范数通过曲线长度产生，因此诱导前一节的局部参数拓扑。Hamenstädt 的综述在第3节陈述了这一结论；此处引用的是作者公开稿的该节，页码与正式文集版本不同。\cite[作者公开稿第3节，第234--235页]{Hamenstadt2006Mapping}
\cref{g4:thm:0602:qc-compactness}给出了这类论证中的基本解析输入：固定伸缩界的规范化映射不能任意坍缩，伸缩比趋于 \(1\) 时其极限恢复共形性。本节后面的长度--扭转参数化将由有限双曲拼接另行证明。

\begin{example}
	\label{g4:exm:0604:annulus-distortion}
	若两个有限圆环域的模量为 \(M_1,M_2\)，则任一准共形同胚的伸缩比至少为
	\[
		\max\{M_2/M_1,M_1/M_2\}.
	\]
	在上述圆柱坐标中，映射
	\((x,\theta)\mapsto(L_2x/L_1,\theta)\)
	达到这个下界，因为它两个方向的伸缩分别为 \(L_2/L_1\) 与 \(1\)。因此环域之间按式\eqref{g4:eq:0604:teich-distance}取下确界时，数值为
	\(\tfrac12|\log(M_2/M_1)|\)。
	这是一个完全可计算的参数模型：很长的圆柱与较短圆柱之间必有明显的非共形伸缩。
\end{example}

\subsection{Fenchel–Nielsen 参数的模型}
\label{g4:subsec:0604:03}

在 \(S\) 上固定一组互不相交、互不同伦且本质的简单闭曲线 \(c_1,\ldots,c_N\)，把 \(S\) 切成三边界球面；每个这样的分块称为一个\term[shuangqukuzi]{双曲裤子}[hyperbolic pair of pants]的拓扑模型。为各块标号，并在每块上选定连接不同边界的三条缝线的同位类。带这些辅助资料的分解将在下面得到几何实现。

\ProseParagraphBreak

每个曲线 \(c_i\) 将给出两个实数：其双曲测地代表的长度 \(\ell_i>0\)，以及把两侧裤子拼回时沿该测地线平移的有向距离 \(\tau_i\in\mathbb R\)。后者的单位是长度；一次完整 Dehn 扭转改变 \(\tau_i\) 一个 \(\ell_i\)。若改用圈数 \(u_i=\tau_i/\ell_i\)，一次完整扭转就改变 \(u_i\) 一个整数。这就是\term[FenchelNielsencanshu]{Fenchel--Nielsen 参数}[Fenchel--Nielsen coordinates]\footnote{Werner Fenchel（1905--1988），德国出生、后在丹麦工作的数学家，研究凸几何与非 Euclid 几何。}的模型：
\[
	(\ell_1,\tau_1,\ldots,\ell_N,\tau_N)
	\in(\mathbb R_{>0}\times\mathbb R)^N.
\]
具体的六边形构造、扭转的提升选择及恢复将在本节后四个小节证明。

\begin{proposition}
	\label{g4:prop:0604:pants-count}
	闭连通定向亏格 \(g\ge2\) 的曲面存在裤子分解；每个裤子分解含 \(2g-2\) 条裤子和 \(N=3g-3\) 条切割曲线。因此长度与扭转总共提供 \(6g-6\) 个实参数。
\end{proposition}
\begin{proof}
	由第一卷命题6.4.5，可以在标准手柄模型中工作。先沿每个手柄的一条非分离曲线切开，得到有 \(2g\) 个边界的球面。每次用一条圈从这个球面分出两个边界与一个新切口组成的裤子，直到余块也是裤子，就得到所需分解。这些切线均本质且不平行；把最初的成对边界重新粘合，便回到 \(S\)。

	由第一卷命题6.5.2，沿边界圆周粘合时 Euler 示性数相加。对任一裤子分解，设裤子数为 \(P\)。每条裤子的 Euler 示性数为 \(-1\)，沿圆周粘合不改变其和，所以 \(-P=2-2g\)，即 \(P=2g-2\)。每条裤子有三条边界，而每条切线贡献两侧，故 \(3P=2N\)，于是 \(N=3g-3\)。同一裤子的两条边界互相拼合时仍计两侧，这个计数同样成立。
\end{proof}

\subsection{退化、节点和紧化的入口}
\label{g4:subsec:0604:04}

参数 \(\ell_i>0\) 描述光滑曲面；令 \(\ell_i\to0\) 会离开这个参数域。在共形坐标中，模型是把两端按 \(zw=t\)、\(t\ne0\) 接成环带，\(t\to0\) 时得到两条分支 \(zw=0\)。第五章已计算其模量
\[
	\operatorname{Mod}=\frac1{2\pi}\log\frac{c^2}{|t|}\longrightarrow\infty
\]
及两侧固定截口观察下的尖点极限。相应的结点不是某个光滑黎曼曲面的普通坐标点，而是两条局部分支相遇的位置。

\ProseParagraphBreak

标记与忘记标记会给出不同的边界行为。固定长度后反复作 Dehn 扭转，可能只改变标记而不改变底层曲面；缩短被扭转的圈则改变了共形几何。这两种操作不能用同一个“参数发散”代替。\secref{g4:sec:0605}将在有序四点球面上通过交比把三个碰撞方向写成具体的结点族；本节则在最后一个小节固定被捏缩的裤子曲线，追踪长度、缝线与扭转。一般紧化还要说明允许哪些奇异曲线、如何比较全部退化族以及标记的遗忘方式。

\subsection{双曲裤子的边界长度与裤子分解}
\label{g4:subsec:0604:05}

双曲裤子的三条边界都是闭测地线。把裤子沿三条最短缝线切开，就应得到两个全等的直角六边形；边界长的一半成为交替的三条边。下面先直接构造这种六边形，借此同时证明任意三个正边界长度都能实现，而且没有隐藏的额外参数。

\begin{lemma}
	\label{g4:lem:0604:right-hexagon}
	给定 \(a,b,c>0\)，存在一个紧凸双曲直角六边形，使三条交替边长为 \(a,b,c\)，并且它在允许反向的等距下唯一。连接 \(a,b\) 两边的边长 \(x\) 满足
	\begin{equation}
		\label{g4:eq:0604:hexagon-cosine}
		\cosh x=\frac{\cosh c+\cosh a\cosh b}{\sinh a\sinh b}.
	\end{equation}
	其余两条边由循环置换得到。
\end{lemma}
\begin{proof}
	采用第一卷定理10.1.3的双曲面距离与测地线公式。在 Lorentz 模型\footnote{Hendrik Antoon Lorentz（1853--1928），荷兰理论物理学家，研究电磁理论、电子理论与运动物体的变换。} \(\langle u,v\rangle=-u_0v_0+u_1v_1+u_2v_2\)、\(\mathbb H^2=\{\langle u,u\rangle=-1,u_0>0\}\) 中，取两条有公垂线的测地线，垂足为
	\[
		p_1=(1,0,0),\qquad p_2=(\cosh x,\sinh x,0),
	\]
	并沿二者相同一侧的单位切向量 \((0,0,1)\) 分别走距离 \(a,b\)，到达
	\[
		q_1=(\cosh a,0,\sinh a),\quad
		q_2=(\cosh b\cosh x,\cosh b\sinh x,\sinh b).
	\]
	在 \(q_1,q_2\) 处作垂直于原两条测地线的新测地线 \(h_1,h_2\)。它们的单位法向量分别为
	\[
		n_1=(\sinh a,0,\cosh a),\quad
		n_2=(\sinh b\cosh x,\sinh b\sinh x,\cosh b).
	\]
	因而 \(\langle n_1,n_2\rangle=\cosh a\cosh b-\sinh a\sinh b\cosh x\)。式\eqref{g4:eq:0604:hexagon-cosine}的右边大于 \(1\)，因为分子减分母为 \(\cosh c+\cosh(a-b)>0\)，所以它唯一确定 \(x>0\)。按此选择有 \(\langle n_1,n_2\rangle=-\cosh c<-1\)。两个单位空间向量张成的平面因此为 Lorentz 平面，正交补为正定直线；在它们张成的平面中取单位双曲向量，直接归约到 \((0,1,0)\) 和 \((\sinh c,-\cosh c,0)\) 的模型，便知两测地线不相交，其唯一公垂线长为 \(c\)。

	还须保证公垂足在所需的正向射线上。沿 \(h_1\) 向第二条原测地线走，参数式为
	\[
		r_1(t)=\cosh t\,q_1+\sinh t\,(0,1,0).
	\]
	令 \(A=-\langle q_1,n_2\rangle=(\cosh a\cosh c+\cosh b)/\sinh a\)，\(B=\sinh b\sinh x\)。则 \(A,B>0\)，且代入上式可算得 \(A^2-B^2=\sinh^2c\)。距离 \(h_2\) 的垂足由 \(\tanh t=B/A\) 确定，所以 \(t>0\)，并有
	\[
		\cosh t=\frac{\cosh b+\cosh a\cosh c}{\sinh a\sinh c}.
	\]
	另一侧完全对称，得到另一个正边长。两足还位于原两测地线之间：例如上面的 \(\tanh t<\cosh a\tanh x\) 正是 \(r_1(t)\) 尚未越过第二条原测地线的条件，代入化为 \(\cosh c<\cosh^2a\cosh c\)。两原测地线围成的带及各垂线的内半平面均测地凸，故这六条边按次序围出一个紧凸六边形，六角都是直角。

	反过来，任一直角六边形可由等距把前述两条边及其间公垂边放到上述位置。相邻垂线的法向内积必为 \(-\cosh c\)，从而强迫式\eqref{g4:eq:0604:hexagon-cosine}；\(x\) 及所有垂线随之确定，其余两条边也唯一。这证明存在与唯一，且给出了全部边长的计算。
\end{proof}

\begin{theorem}
	\label{g4:thm:0604:pants-lengths}
	给定三个正数 \(\ell_1,\ell_2,\ell_3\)，存在边界依次标号的紧双曲裤子，其测地边界长为这些数；它在保向且保持边界标号的等距下唯一。连接第 \(i,j\) 条边界的最短缝线长 \(s_{ij}\) 满足
	\[
		\cosh s_{ij}
		=\frac{\cosh(\ell_k/2)+\cosh(\ell_i/2)\cosh(\ell_j/2)}
		{\sinh(\ell_i/2)\sinh(\ell_j/2)},\qquad \{i,j,k\}=\{1,2,3\}.
	\]
\end{theorem}
\begin{proof}
	将两个交替边为 \(\ell_1/2,\ell_2/2,\ell_3/2\) 的六边形沿另外三边反射粘合。测地边的反射是局部双曲等距，内点两侧的度量因此光滑；在各端点，两直角合成平角，所以留下的边界也是光滑测地线。拓扑上得到三边界球面，且三条边界长正好为 \(\ell_1,\ell_2,\ell_3\)。

	反向考虑任一双曲裤子。不同边界之间的距离在紧性下达到；最短弧的内部是测地线，端点自由变分使它与两边界垂直。它不能在途中接触第三条测地边界，否则相切后由测地线唯一性不能再离开。三对边界的最短弧彼此不交：若两弧从同一个边界到达一个交点，交换从该边界到交点的两段，较短的一段会缩短另一条弧；若两段等长，交换后在交点形成的角又能严格缩短。重叠一段后再分开也违反测地线初值唯一性。

	三条不交弧把裤子切成两个圆盘，每个圆盘的边界有六条边且六角为直角。这个拓扑事实可在球面去三个开盘的模型中沿弧逐条切开核对。两六边形的三条交替缝线边分别相等，故由上一引理唯一性，两六边形全等。于是各边界被分成相等两半，边长恰为 \(\ell_i/2\)。引理便同时给出唯一性与缝线公式；反射对称也说明上述三条缝线由几何唯一确定。
\end{proof}

% !TeX root = ../Book4.tex
\begin{figure}[htbp]
	\centering
	\begin{tikzpicture}[font=\small,scale=1.05]
		\fill[GeometryBlueLight] (0,0) ellipse (3 and 1.65);
		\fill[white] (-1.2,0) circle (0.62);
		\fill[white] (1.2,0) circle (0.62);
		\draw[GeometryGray,thick] (0,0) ellipse (3 and 1.65);
		\draw[GeometryGray,thick] (-1.2,0) circle (0.62);
		\draw[GeometryGray,thick] (1.2,0) circle (0.62);
		\draw[GeometryBlue,very thick] (-3,0)--(-1.82,0);
		\draw[GeometryBlue,very thick] (-0.58,0)--(0.58,0);
		\draw[GeometryBlue,very thick] (1.82,0)--(3,0);
		\node[above] at (-2.4,0.05) {\(s_{13}\)};
		\node[above] at (0,0.05) {\(s_{12}\)};
		\node[above] at (2.4,0.05) {\(s_{23}\)};
		\node at (-1.2,0) {\(C_1\)};
		\node at (1.2,0) {\(C_2\)};
		\node[above] at (0,1.65) {\(C_3\)};
		\node at (0,1) {上侧六边形};
		\node at (0,-1) {下侧六边形};
		\node[below] at (-1.2,-0.62) {\(\ell_1/2\)};
		\node[below] at (1.2,-0.62) {\(\ell_2/2\)};
		\node[below] at (0,-1.65) {\(\ell_3/2\)};
	\end{tikzpicture}
	\caption[裤子的三条缝线与六边形分解]{三条蓝色缝线分别连接不同边界，把裤子分成上下两个直角六边形。每个六边形的三条边界弧长为 \(\ell_1/2,\ell_2/2,\ell_3/2\)，另三边为缝线。图为平面拓扑示意，所画圆周和长度比例不表示双曲测量。}
	\label{g4:fig:0604:pants-seams}
\end{figure}


\cref{g4:fig:0604:pants-seams}把三条缝线与两个六边形的对应画出；图中的圆周只表示边界的拓扑位置，六角的直角性来自双曲度量。

对闭双曲曲面，所选裤子曲线都能同时换成互不相交的测地线。其原因可以从万有覆盖看清：一个非平凡自由同伦类对应一个覆盖变换。紧商可由有限个均匀覆盖的小圆盘覆盖，取这些盘的 Lebesgue 数便给出非平凡覆盖变换位移的正下界。由第一卷定理10.2.1，该变换不能是抛物型；自由作用又排除椭圆型，所以它有唯一双曲轴，投影得到该类的闭测地线。原简单曲线的一条提升由一个紧基本弧在该覆盖变换下的像拼成，因而与双曲轴保持有界距离；它是连接轴的两个理想端点的嵌入直线。不同提升不相交。若相应轴相交，其端点在理想圆周上交错，第一卷定理6.2.3的平面分离迫使原提升也相交，矛盾。因此测地代表仍简单；对两条不同且不交的曲线作同样比较，测地代表也不交。由\cref{g4:prop:0601:surface-isotopy-interface}，简单闭曲线的自由同伦与同位相容，所以这些测地线代表原来指定的切割同位类。切开以后正是上一定理描述的双曲裤子。

\subsection{长度、扭转与标记的选择}
\label{g4:subsec:0604:06}

边界长为 \(\ell\) 的测地圆周，可用弧长坐标 \(t\in\mathbb R/\ell\mathbb Z\) 表示。给每条裤子边界指定其中一条缝线的垂足作为零点，并固定与参考曲线方向相容的坐标方向。两侧边界要反向粘合以保持曲面定向；若各用本侧诱导的边界方向，拼接式写为
\[
	t_-=-t_++\tau\pmod\ell.
\]
若两侧改用同一个参考曲线方向，则它就是一个平移。几何上的拼接只看见 \(\tau\bmod\ell\)，而标记还记录绕行的整数次数。

\ProseParagraphBreak

为把这个整数固定下来，在参考曲面的每个切割环带中选一条横穿弧，连接两侧指定的缝线。将标记送来的这条弧与几何缝线对齐端点，在环带的条带覆盖中提升；规定左端弧长坐标为零后，右端坐标的实数值就是 \(\tau\)。更换弧的相对同伦类一个正整圈，会使 \(\tau\) 增加 \(\ell\)。所以 \(\tau\in\mathbb R\) 由曲面、标记及固定的缝线与横穿弧共同确定，不能在保留标记时把它任意约到 \([0,\ell)\)。

\ProseParagraphBreak

Dehn 扭转的符号可按这一条带坐标固定：称使提升右端增加 \(\ell\) 的扭转为本参数的正扭转；其对标记的作用为 \((\ell,\tau)\mapsto(\ell,\tau+\ell)\)。若映射类作用使用预合成逆映射的约定，同一个群元素在参数上的符号也相应反转。零点换用另一条缝线通常增加半个边界长度，改变横穿弧则增加整倍边界长度；这些都是辅助选择的改变，不是曲面新添了连续参数。

\subsection{沿测地边界的拼接与参数恢复}
\label{g4:subsec:0604:07}

\begin{theorem}
	\label{g4:thm:0604:fenchel-nielsen}
	固定亏格 \(g\ge2\) 的闭连通定向参考曲面 \(S\)，以及上述带标号的裤子分解、缝线和横穿弧，令 \(N=3g-3\)。长度与实扭转给出一一对应
	\[
		\mathcal T(S)\longrightarrow
		(\mathbb R_{>0}\times\mathbb R)^N,
		\qquad [X,f]\longmapsto(\ell_i,\tau_i)_{i=1}^N.
	\]
	在局部 Teichmüller 拓扑与右边的通常拓扑下，它及其逆均连续。
\end{theorem}
\begin{proof}
	先从任意一组参数构造曲面。每条裤子的三条边界长由相邻 \(\ell_i\) 指定，\cref{g4:thm:0604:pants-lengths}给出它的唯一模型。沿所有配对边界按 \(t_-=-t_++\tau_i\) 粘合。每条测地边界附近有法向坐标，度量为
	\[
		\mathrm{d}r^2+\cosh^2r\,\mathrm{d}t^2,
	\]
	其中边界为 \(r=0\)。该式可在上半平面轴旁直接计算，正是\cref{g4:eq:0504:cylinder-metric}的局部版本。反向配对以后以有符号的 \(r\) 合并两侧，\(\cosh r\) 为偶函数，度量光滑通过接缝，且曲率仍为 \(-1\)。有限块粘合成紧定向曲面；其全纯结构由局部双曲图给出。

	标记在每条裤子上由缝线和边界标号确定，再在各环带中用所指定实位移的条带提升接起来。相差整周的位移给同一个未标记的边界配对，却使横穿弧多绕一圈，因而按定义改变标记的 Dehn 扭转数。这说明全部实数 \(\tau_i\) 都有实现。

	反向给定标记曲面，以统一化的 Poincaré 度量将各 \(c_i\) 化为测地代表并切开。边界长给出 \(\ell_i\)，每条裤子的唯一模型确定其缝线位置；标记的横穿弧在条带覆盖中的终点给出 \(\tau_i\)。因此构造后再读取参数会恢复原数据。若两张标记曲面参数相同，各裤子的保向等距在缝线、边界及横穿弧上相容，遂拼成全局等距，其标记差固定所选切割弧；在剩下的圆盘中以及扭转数为零的环带中，\cref{g4:lem:0601:relative-disk-annulus}将它相对同位为恒等。它们代表同一个 Teichmüller 点。等距是双全纯的；反过来双全纯等价保持 Poincaré 度量与标记数据，所以对应在定义的等价类上成立。

	最后核对连续性。六边形的坐标、边长和垂足公式在全部 \(\ell_i>0\) 的区域光滑。用缝线、法向领圈及内部有限三角片识别邻近裤子，得到随参数连续且局部光滑的标记度量；领圈中的扭转由一个固定截断函数乘 \(\tau_i/\ell_i\) 实现。因此逆构造在局部切片拓扑中连续。反向在局部切片中选取连续依赖参数的光滑双曲度量。各固定类的测地代表长度局部有界；按弧长参数化后，由紧性与测地线方程，其任一序列都有光滑收敛子列。极限仍在原自由同伦类，并为测地线，唯一性使它就是极限度量的代表；因此全序列收敛。裤子内垂直缝线同样由最短性、紧性和唯一性连续变化。由此长度和扭转模 \(\ell_i\) 连续；横穿弧固定的整数提升在足够小邻域中不跳变，所以实扭转也连续。
\end{proof}

这个参数化的实维数为 \(2(3g-3)\)，与\cref{g4:cor:0603:dimension}的复维数 \(3g-3\) 相符。长度和扭转是实几何坐标，不能逐对把 \(\ell_i+\mathrm{i}\tau_i\) 直接当作前一节的全纯坐标；两种坐标分别服务于不同的计算。

\subsection{长度趋零时的标记与退化}
\label{g4:subsec:0604:08}

固定裤子分解中的一条曲线 \(c\)，让它的长度 \(\ell\downarrow0\)，并把其余切割曲线的长度限制在正的紧区间中。缝线公式立刻显示通向这条边界的缝线越来越长。例如，若邻近裤子的另外两条边界长 \(b,c_0\) 保持正值，则
\[
	\cosh s
	=\frac{\cosh(c_0/2)+\cosh(\ell/2)\cosh(b/2)}
	{\sinh(\ell/2)\sinh(b/2)},
\]
所以 \(s=\log(1/\ell)+O(1)\)。若切线两侧来自同一条裤子，便有两条边界同时长为 \(\ell\)：连接这两条边界的缝线长为 \(2\log(1/\ell)+O(1)\)，它们到第三条正长度边界的缝线仍为 \(\log(1/\ell)+O(1)\)，直接代入同一公式即可得到。原来有限长度的边界因而远离固定的内部部分；将未收缩部分固定在模型中，相关六边形顶点趋于理想边界，得到以一个或两个尖点代替相应测地边界的裤子。

\ProseParagraphBreak

对指定的嵌入双曲环带，这个变化有更精确的尺度。\cref{g4:prop:0505:truncated-cylinder}计算了核心长 \(\ell\)、两侧截口长固定为 \(a>\ell\) 的对称模型：其模量为
\[
	M_\ell=\frac2\ell\arccos(\ell/a)
	=\frac\pi\ell-\frac2a+O(\ell^2).
\]
从任一截口向内量距离 \(s\)，度量系数为
\[
	\mathrm{d}s^2+
	\bigl(a\cosh s-\sqrt{a^2-\ell^2}\sinh s\bigr)^2\,\mathrm{d}x^2,
	\qquad x\in\mathbb R/\mathbb Z.
\]
它在每个固定宽度区域上光滑趋于 \(\mathrm{d}s^2+a^2e^{-2s}\mathrm{d}x^2\)。这精确表达了两侧看到两个尖点，而中间核心退到无限距离处的现象。对当前裤子拼接族，可以在理想六边形极限的尖点内先选足够短的截口，再取邻近参数的对应截口，得到这样的嵌入端部；这里处理的是指定拼接族，未对所有短闭测地线宣称同一个全局领圈宽度。

\ProseParagraphBreak

标记仍记住这条被捏缩的曲线及横穿它的弧。实扭转 \(\tau\) 与圈数 \(u=\tau/\ell\) 的行为须分别记录：即使 \(\tau\) 有界，\(u\) 也可能无界；而作一次 Dehn 扭转改变 \(\tau\) 的量本身趋于零。忘记标记以后，\(u\) 的整数部分被相应映射类识别。将两个尖点补成结点时，原来的边界圆不再存在，沿该圆的扭转也不再作为结点曲线本身的一个参数保留；邻近光滑族的粘合参数 \(t\) 仍通过其辐角记录相对旋转。

\ProseParagraphBreak

若同时捏缩若干条互不相交的裤子曲线，可在各配对边界分别进行上述构造。切开后每个分块具有负 Euler 示性数；把尖点按原配对补为结点，所得分支满足稳定性条件：球面分支至少有三个特殊点，环面分支至少有一个。其余正长度和扭转参数描述未捏缩部分。这个有限模型解释了光滑参数域怎样邻接结点曲线，并为下一节四点球面的三个具体边界提供度量图像。
% 纲要条目排为编号小节标题，并进入目录与书签；说明段仍保留于纲要注释中。

% !TeX root = ../../Book4.tex
% 纲要标题：### 6.5 四点标记球面的参数与退化
% 纲要位置：几何学计划纲要.md，第 3726 行。
% 正文状态：初稿；尚未正式审核。

% 纲要内容开始
% 
% 这里完成解析参数和三类退化的识别；跨任意代数基底的通用性证明在 V.34.6 进行。交比参数化的是模空间，不把它等同于其带标记的 Teichmüller 覆盖。
% 
% * 有序四点、Möbius 归一化与交比参数
% * $\mathbb P^1\setminus\{0,1,\infty\}$ 上的光滑四点族
% * 标记碰撞、长颈与结点极限
% * 两两分配标记的三种稳定边界类型
% * 有序标记与置换商、自同构的区别
% * 与 V.34.6 的代数通用族及完整紧化构造比较
% 
% ---
% 
% 纲要内容结束

\section{四点标记球面的参数与退化}
\label{g4:sec:0605}

球面本身没有复模参数，但在球面上指定四个有序的点后，第四个点相对于前三点的位置不能再由坐标变化消去。这个一维问题可以直接算完：光滑对象由一个交比描述，参数逼近三个缺失值时，碰撞的点分别落到新出现的球面分支上。下面同时保留参数、标记和局部方程，以便看清模空间中的边界点究竟表示什么。

\subsection{有序四点、Möbius 归一化与交比参数}
\label{g4:subsec:0605:01}
设 \(p_1,p_2,p_3,p_4\) 是 \(\mathbb P^1\) 上两两不同的点。这里的有序标记是四个点的编号；本章前面的一张拓扑曲面到复曲面的标记则包含道路和同位信息。四点的编号本身不记录这些道路。

\begin{proposition}
	\label{g4:prop:0605:cross-ratio}
	球面的双全纯自同构恰为 \(z\mapsto(az+b)/(cz+d)\)、\(ad-bc\ne0\)。任意三个不同的有序点有唯一的这类变换 \(N\)，满足 \(N(p_1)=0\)、\(N(p_2)=1\)、\(N(p_3)=\infty\)。数值
	\[
		\lambda=N(p_4)=\frac{(p_4-p_1)(p_2-p_3)}{(p_4-p_3)(p_2-p_1)}
	\]
	按无穷远处的极限解释，是有序四点球面在保持编号的双全纯同构下的完全不变量，其取值集为 \(B=\mathbb P^1\setminus\{0,1,\infty\}\)。
\end{proposition}
\begin{proof}
	先解释自同构的形式。设 \(F\) 为球面自同构，把唯一的 \(F^{-1}(\infty)\) 用一个分式线性坐标移到无穷远。此后 \(F\) 在有限平面全纯，并因无穷远处局部可逆而只有一阶极点。因此 \(|F(z)|\le C(1+|z|)\)，Cauchy 估计使二阶导数恒为零，\(F\) 是斜率非零的仿射函数。换回坐标得到所述分式形式。反之矩阵行列式不为零时，逆矩阵给出全纯逆。

	当前三点有限时，可取 \(N(z)=(z-p_1)(p_2-p_3)/((z-p_3)(p_2-p_1))\)；有无穷远点时取极限。若两变换均满足归一化，其复合逆固定 \(0,1,\infty\)：固定无穷远迫使它仿射，固定零和一再使它恒等。故归一化唯一。第四点不同于前三点等价于 \(\lambda\ne0,1,\infty\)。两组四点交比相同，当且仅当各自归一化后都是 \((0,1,\infty,\lambda)\)，因而可由保持编号的双全纯映射联系。
\end{proof}

这类分式变换称为 Möbius 变换。若将四个点看成球面配置空间的点，上述归一化在各坐标片上是全纯的，且分母只在被排除的碰撞处为零。因此交比不仅给出集合分类，也给出这个配置空间商的一维复坐标。我们记有序四点球面的模空间为 \(\mathcal M_{0,4}\)，便有 \(\mathcal M_{0,4}\simeq B\)。

\subsection{\texorpdfstring{\(\mathbb P^1\setminus\{0,1,\infty\}\) 上的光滑四点族}{P¹∖\{0,1,∞\} 上的光滑四点族}}
\label{g4:subsec:0605:02}
在交比坐标下，所有光滑对象同时出现在一个简单的族中：
\[
	\pi:\mathcal C=\mathbb P^1\times B\longrightarrow B,
	\qquad s_1(\lambda)=0,\quad s_2(\lambda)=1,\quad
	s_3(\lambda)=\infty,\quad s_4(\lambda)=\lambda.
\]
这里 \(s_i\) 是投影的截面，简写中省去了共同的第二坐标 \(\lambda\)。投影是固有的全纯下沉，四个截面处处不交。给定 \(\lambda\) 的纤维就是前一命题的标准代表。

\begin{proposition}
	\label{g4:prop:0605:family-normalization}
	设 \(T\) 为复流形，\(C\to T\) 是局部全纯平凡的球面族，带四个两两不交的有序全纯截面。则存在唯一全纯映射 \(\lambda:T\to B\)，以及唯一保持四个截面的族同构
	\[
		C\simeq\lambda^*\mathcal C,
	\]
	右边表示按 \(\lambda\) 拉回上述标准族。
\end{proposition}
\begin{proof}
	在一个平凡化邻域中，前三个截面给出三个随底点全纯变化且不相等的球面点。用前一命题的公式归一化；必要时先换一张球面坐标，使所用分母在当前底点非零，再缩小邻域。由公式可见归一化分式的系数随底点全纯变化，故给出局部族同构。第四截面在其中的坐标就是全纯函数 \(\lambda\)。交叠处的两个归一化都把前三截面送到 \(0,1,\infty\)，由唯一性逐纤维相等，因而作为全纯映射相等。它们遂粘合为全局同构及全局参数。由此得到交换图
	\[
		\begin{tikzcd}
			C \arrow[r] \arrow[d] & \mathcal C \arrow[d,"\pi"] \\
			T \arrow[r,"\lambda"'] & B.
		\end{tikzcd}
	\]
	上箭头在每条纤维上保持编号。唯一性同样归结为固定三个点的自同构恒等。
\end{proof}

\(B\) 不是单连通的：绕 \(0\)、\(1\) 的道路保留非平凡绕行信息。给球面四个点再指定拓扑标记，沿 \(B\) 中闭路移动第四点可能改变这个标记。前面定义的带标记空间记录这类信息；交比记录的是遗忘拓扑标记后的对象。上述族的构造和分类在整个 \(B\) 上成立，不需要选择一个交比的对数或一条回到基点的道路。

\subsection{标记碰撞、长颈与结点极限}
\label{g4:subsec:0605:03}
先令 \(\lambda=t\to0\)。在固定的 \(z\) 坐标中，第1和第4标记碰撞于零，极限不再有四个不同标记。为了保留两点的相对位置，引入放大坐标 \(\xi=z/t\)。这时两点始终为 \(\xi=0,1\)，而原先留在远处的部分从 \(\xi=\infty\) 接入。一个坐标适合远处，另一个适合碰撞的近处，它们需要在同一个曲面族中兼容。

\begin{proposition}
	\label{g4:prop:0605:node-family}
	取 \(|t|<\varepsilon<1/2\)。将 \(\mathbb P^1_z\times\{|t|<\varepsilon\}\) 的点 \((0,0)\) 替换为记录方向的射影直线，局部用
	\[
		zV=tU,\qquad [U:V]\in\mathbb P^1
	\]
	表示。所得复曲面光滑，投影到 \(t\) 为固有映射。\(t\ne0\) 的纤维仍是原四点球面；\(t=0\) 的纤维是两个球面在一个结点相交，每个球面恰含两个标记及该结点。结点附近的族方程为 \(xy=t\)。
\end{proposition}
\begin{proof}
	在 \(V\ne0\) 的坐标中令 \(\xi=U/V\)，方程为 \(z=t\xi\)，总空间以 \((t,\xi)\) 为光滑坐标。在 \(U\ne0\) 的坐标中令 \(\eta=V/U\)，方程为 \(t=z\eta\)，总空间以 \((z,\eta)\) 为光滑坐标。交叠处 \(\eta=1/\xi\)、\(z=t\xi\)，它们与未修改区域中的原坐标全纯相容。投影到原乘积是闭的射影方向条件，且方向因子紧，故为固有映射；再复合原来的固有投影即得所需固有性。对 \((z,t)\ne(0,0)\)，方向唯一，故纤维在 \(t\ne0\) 时不变。

	在第二张图中，中心纤维为 \(z\eta=0\)。\(z=0\) 是新增的方向球面，\(\eta=0\) 是原中心球面去掉零点后的闭包，两者在 \((z,\eta)=(0,0)\) 相交。第一张图中的 \(t=0\) 补全新增球面，其 \(\xi=\infty\) 就是这个交点。因此结点方程为 \(xy=t\)，取 \(x=z\)、\(y=\eta\)。

	第1、第4截面在第一张图中分别为 \(\xi=0\)、\(\xi=1\)，中心点各不相同，也不在结点。第2、第3截面仍位于原球面上的 \(z=1,\infty\)。于是各分支均有三个不同的特殊点，即两个标记和一个结点。任何保持所有标记的自同构必须保持两个分支，并在每个分支固定三个点，故为恒等。
\end{proof}

在结点图中固定 \(0<c<1\)，当 \(0<|t|<c^2\) 时，\(|x|<c\)、\(|y|<c\)、\(xy=t\) 的部分等价于
\[
	|t|/c<|x|<c.
\]
由\eqref{g4:eq:0505:plumbing-modulus}，其模量为 \(\log(c^2/|t|)/(2\pi)\)。所以碰撞不仅产生一个形式上的奇点，也产生模量趋于无穷的实际环形颈部。若 \(t=|t|e^{\mathrm{i}\theta}\)，粘合关系给出 \(\arg y=\theta-\arg x\)；辐角控制两端的相对转动，而模控制颈部的长度尺度。

\subsection{两两分配标记的三种稳定边界类型}
\label{g4:subsec:0605:04}
对于亏格零的结点曲线，稳定性的要求是：每个球面分支上的特殊点总数至少为三；特殊点包括标记点以及结点在该分支上的原像。这个条件正好排除固定一两个点仍可连续变化的球面自同构。

\ProseParagraphBreak

结点曲线的对偶图以归一化分支为顶点、结点为边。其算术亏格在这个模型中为各分支亏格之和加上对偶图的圈数。这个数也等于把各结点换成小颈部后光滑曲面的亏格：每挖去一个结点原像的小盘，Euler 示性数减少1；沿圆周接入圆柱不再改变它。对有 \(v\) 个分支、\(e\) 个结点的连通曲线，所得等式为 \(g=\sum g_i+e-v+1\)。所以算术亏格零恰要求所有分支为球面、对偶图为树。

\begin{proposition}
	\label{g4:prop:0605:three-boundaries}
	连通、算术亏格零、带四个有序光滑标记的稳定结点曲线，若非光滑，则恰由两个球面组成，各带两个标记，并在一个结点相接。其同构类型恰有三种，分别由无序二二分割
	\[
		\{1,4\}\mid\{2,3\},\qquad
		\{2,4\}\mid\{1,3\},\qquad
		\{3,4\}\mid\{1,2\}
	\]
	确定，它们分别对应 \(\lambda=0,1,\infty\)。
\end{proposition}
\begin{proof}
	亏格零意味着归一化的分支都是球面，连接各分支的对偶图是树。设有 \(v\) 个分支，则有 \(v-1\) 个结点。特殊点在各分支上的总数为 \(4+2(v-1)=2v+2\)。每个分支至少有三个，故 \(3v\le2v+2\)，即 \(v\le2\)。非光滑时 \(v=2\)，两边各有一个结点原像，稳定性遂迫使各自恰有两个标记。每个分支的三个特殊点都可唯一归一化为 \(0,1,\infty\)，所以一旦指定二二分割，曲线的同构类及其保持标记的同构就都确定。

	\(\lambda=0\) 的分割已由前一命题得到。\(\lambda\to1\) 时取 \(t=\lambda-1\)、\(x=z-1\)，碰撞的是第2和第4标记；重复同一坐标构造得到第二种。\(\lambda\to\infty\) 时取 \(t=1/\lambda\)、\(x=1/z\)，碰撞的是第3和第4标记，得到第三种。三个局部构造在非零参数处都恢复原族，因此可与 \(B\) 上的标准族粘合，给出参数球面上的稳定解析族。
\end{proof}

% !TeX root = ../Book4.tex
\begin{figure}[htbp]
	\centering
	\begin{tikzpicture}[scale=0.9, font=\small]
		\foreach \x/\a/\b/\c/\d/\v in {0/1/4/2/3/0,4.5/2/4/1/3/1,9/3/4/1/2/\infty}{
			\begin{scope}[xshift=\x cm]
				\draw[thick,GeometryBlue] (0,0)--(1.4,0);
				\draw (0,0)--(-0.6,0.7) node[above]{\(\a\)};
				\draw (0,0)--(-0.6,-0.7) node[below]{\(\b\)};
				\draw (1.4,0)--(2,0.7) node[above]{\(\c\)};
				\draw (1.4,0)--(2,-0.7) node[below]{\(\d\)};
				\fill[GeometryBlue] (0,0) circle (2pt) (1.4,0) circle (2pt);
				\node[below] at (0.7,-1.2){\(\lambda=\v\)};
			\end{scope}
		}
	\end{tikzpicture}
	\caption[有序四点球面的三个稳定边界]{三个边界纤维的对偶图。每个球面分支有两个标记和一个结点原像，三种分割由标记的编号区分。}
	\label{g4:fig:0605:boundary}
\end{figure}


图中的顶点代表球面分支，连接边代表结点，短枝的编号代表标记。它显示的是对偶图，不是将复曲线嵌入平面的图像。每个顶点恰连着三条枝或边，与稳定性条件相符。

\subsection{有序标记与置换商、自同构的区别}
\label{g4:subsec:0605:05}
保留四点的次序时，每个光滑对象的自同构群都平凡；若只记四点集合，自同构可以置换这些点。忘掉编号既改变参数的等价关系，也改变对象的对称群。

\begin{proposition}
	\label{g4:prop:0605:permutations}
	重排四个标记产生的交比恰属于
	\[
		\lambda,\quad 1-\lambda,\quad\lambda^{-1},\quad
		(1-\lambda)^{-1},\quad\frac{\lambda}{\lambda-1},\quad
		\frac{\lambda-1}{\lambda}.
	\]
	由 \(S_4\) 到这六个分式变换组成的群的同态，核为四元群
	\(V_4=\{1,(12)(34),(13)(24),(14)(23)\}\)。无序四点球面的自同构群通常有4个元素；当交比属于 \(\{-1,1/2,2\}\) 时有8个元素，当 \(\lambda^2-\lambda+1=0\) 时有12个元素。
\end{proposition}
\begin{proof}
	交换第1与第2点，把标准配置重新归一化，得到 \(1-\lambda\)；交换第2与第3点得到 \(\lambda/(\lambda-1)\)。代入公式还得到逆数变换。这些变换相互复合只有所列六个有理函数，并在三个缺失值上实现全部置换，所以它们组成 \(S_3\)。逐项代入三个双换位，交比均不变。群同态满射且核有 \(24/6=4\) 个元素，故核正是 \(V_4\)。

	固定一个配置后，能由 Möbius 变换实现的点置换，恰是保持该配置交比的置换；存在性来自归一化，唯一性来自前三点。其数目等于 \(4\) 乘以 \(\lambda\) 在上述 \(S_3\) 作用下的稳定子阶数。三个非平凡对合的固定点在 \(B\) 内分别为 \(1/2,-1,2\)，两个三阶变换的固定点满足 \(\lambda^2-\lambda+1=0\)。两类集合不交，故稳定子阶数分别为 \(2,3\)，其他点为 \(1\)，得到所述计数。
\end{proof}

例如交比为 \(-1\) 的配置具有额外对称，因而无序参数的商映射在那里出现二重分歧。它仍是光滑四点球面，绝不是结点边界。有限对称造成的商分歧，与 \(\lambda\to0,1,\infty\) 时产生新分支的退化，是两种不同的几何现象。

\subsection{与 V.34.6 的代数通用族及完整紧化构造比较}
\label{g4:subsec:0605:06}
上述计算得到一个完整的一维解析模型：\(B\) 参数化光滑有序四点球面；补上三个参数值，得到三种稳定结点纤维。局部 \(xy=t\) 将参数的消失、环域模量发散与核心圈收缩联系起来。保持标记的自同构在光滑纤维及三个边界纤维上均平凡，因此这些具体对象没有因对称性产生的识别歧义。

\ProseParagraphBreak

第五卷第34.6节将把同一模型放在代数基底上：从 \(\mathbb P^1_z\times\mathbb P^1_\lambda\) 的四个截面出发，在三个碰撞中心作吹起，并验证平坦性、基变换与通用性。本节的方向坐标已经给出了复数域上相应的局部方程。代数的通用性还要求对非约化测试基底构造拉回及下降，因此除了几何点分类，还需检查无穷小底空间上的族；这些是第五卷的具体证明任务。
% 纲要条目排为编号小节标题，并进入目录与书签；说明段仍保留于纲要注释中。


\begin{chaptersummary}
拓扑标记让同一条道路在形变过程中可以追踪。映射类群通过改变标记作用于参数空间，忘掉标记得到模空间；稳定子记录复曲面的自同构。Dehn 扭转在环带中可见，在同调上表现为保持交叉配对的整数剪切，而 Torelli 核记录同调看不见的映射类。

\ProseParagraphBreak

Beltrami 微分把复结构的变化写为一阶方程。向量场产生的变化对应 \(\bar\partial\) 的像，商空间为 \(H^1(X,T_X)\)，其对偶是全纯二次微分空间。亏格至少二时，两者的复维数都是 \(3g-3\)；实际局部参数把这些线性方向实现为复结构的变化。长度与扭转参数则给出相应的 \(6g-6\) 个实参数，标记保留每条粘合曲线上的整数扭转。

\ProseParagraphBreak

交比将有序四点球面的模空间写成三点穿孔球面。补回三个缺失参数值时，四个标记分别按三种二二分割落到两个球面上；每个分支有三个特殊点，保持编号的自同构因而平凡。结点方程 \(xy=t\) 把这个代数形状与模量趋于无穷的解析长颈联系起来。
\end{chaptersummary}

\BookExercises
\begin{exercise}
	\label{g4:ex:0602:constant-torus}
	设 \(\operatorname{Im}\tau>0\)、\(|\mu|<1\)。在 \(\mathbb C/(\mathbb Z+\tau\mathbb Z)\) 上以常系数方程 \(w_{\bar z}=\mu w_z\) 定义新复结构。证明归一化后的标记参数为
	\[
		\tau_\mu=\frac{\tau+\mu\bar\tau}{1+\mu},\qquad
		\operatorname{Im}\tau_\mu=\frac{(1-|\mu|^2)\operatorname{Im}\tau}{|1+\mu|^2}.
	\]
	求从给定 \(\tau'\in\mathbb H\) 恢复 \(\mu\) 的公式，并证明相应的仿射映射在同一标记类的保向光滑微分同胚中具有最小伸缩比。
\end{exercise}
\begin{solution}
	取 \(w=z+\mu\bar z\)。其正实 Jacobian 为 \(1-|\mu|^2\)，故为保向实线性同构，并满足方程。两条标记周期分别变为 \(1+\mu\)、\(\tau+\mu\bar\tau\)，再除以第一周期即得 \(\tau_\mu\)。与共轭相减并通分，分子为 \((1-|\mu|^2)(\tau-\bar\tau)\)，得到虚部公式。解线性方程得
	\[
		\mu=\frac{\tau-\tau'}{\tau'-\bar\tau}.
	\]
	分母非零，且 \(|\tau'-\bar\tau|^2-|\tau'-\tau|^2=4\operatorname{Im}\tau'\operatorname{Im}\tau>0\)，所以所得 \(|\mu|<1\)。这里的参数改变保留了原来两条道路的顺序；再作整数换基才进入无标记环面的同构问题。

	设保留标记的另一映射在平面上的提升为 \(F(z)=az+b\bar z+u(z)\)，其中 \(u\) 对原格周期，\(a,b\) 由两条目标周期唯一确定，且 \(b/a=\mu\)。周期函数的一阶导数在基本平行四边形上的平均为零；这是散度公式与相对边的抵消。因此 \(\int|F_z|^2\ge A|a|^2\)，其中 \(A=\operatorname{Im}\tau\)。面积换元又给出 \(\int J_F=A(|a|^2-|b|^2)\)，因为提升在商上为一次映射，而右端就是仿射映射的目标面积。设 \(k=\sup|F_{\bar z}/F_z|\)，则
	\[
		A(|a|^2-|b|^2)=\int(|F_z|^2-|F_{\bar z}|^2)
		\ge(1-k^2)\int|F_z|^2\ge(1-k^2)A|a|^2.
	\]
	故 \(k\ge|b/a|=|\mu|\)，相应 \(K(F)\ge(1+|\mu|)/(1-|\mu|)\)。仿射映射处处达到这一系数，故实现所需下界。
\end{solution}

\begin{exercise}
	设 \(f,g\) 为平面区域间保向 \(C^1\) 微分同胚，\(\mu_f=f_{\bar z}/f_z\)、\(\mu_g=g_{\bar w}/g_w\)。推导 \(g\circ f\) 的 Beltrami 系数，并由此证明两张局部可逆的同系数解之转移映射全纯。
\end{exercise}
\begin{solution}
	实链式法则给出
	\[
		(g\circ f)_z=(g_w\circ f)f_z+(g_{\bar w}\circ f)\overline{f_{\bar z}},\qquad
		(g\circ f)_{\bar z}=(g_w\circ f)f_{\bar z}+(g_{\bar w}\circ f)\overline{f_z}.
	\]
	设 \(a=\overline{f_z}/f_z\)，\(|a|=1\)，便有
	\[
		\mu_{g\circ f}=\frac{\mu_f+a(\mu_g\circ f)}{1+a\bar\mu_f(\mu_g\circ f)}.
	\]
	保向性给出各系数模小于一，分母因而非零。若 \(h=g\circ f\) 与 \(f\) 同系数，交叉相乘消去 \(\mu_f\)，得到 \(a(\mu_g\circ f)(1-|\mu_f|^2)=0\)，故 \(\mu_g=0\)。对两张同系数解 \(f,h\)，取 \(g=h\circ f^{-1}\)，即得转移映射全纯。此计算也解释了为什么仅给出系数还需检查图册的相容性。
\end{solution}

\begin{exercise}
	在圆柱 \(C_j=[0,L_j]\times(\mathbb R/2\pi\mathbb Z)\)、\(L_j>0\) 上取平直度量 \(\mathrm{d}x^2+\mathrm{d}\theta^2\)，并固定边界标记。给定 \(T\in\mathbb R\)，考虑保向 \(C^1\) 微分同胚，其提升满足
	\[
		\widetilde F(0,\theta)=(0,\theta),\qquad
		\widetilde F(L_1,\theta)=(L_2,\theta+T).
	\]
	令 \(\alpha=L_2/L_1\)、\(\beta=T/L_1\) 及 \(S=(\alpha^2+\beta^2+1)/\alpha\)。证明在这一边界标记类中，最小伸缩比为 \((S+\sqrt{S^2-4})/2\)，并写出达到它的映射。解释整数圈扭转为何在此问题中仍可辨认。
\end{exercise}
\begin{solution}
	仿射映射 \(F_0(x,\theta)=(\alpha x,\theta+\beta x)\) 满足所需条件。其微分矩阵为
	\[
		A=\begin{pmatrix}\alpha&0\\\beta&1\end{pmatrix}.
	\]
	若奇异值为 \(s_1\ge s_2>0\)，则 \(s_1s_2=\alpha\)、\(s_1^2+s_2^2=\alpha^2+\beta^2+1\)。故其伸缩比 \(K_0=s_1/s_2\) 满足 \(K_0+K_0^{-1}=S\)，于是 \(K_0=(S+\sqrt{S^2-4})/2\)。

	对任一允许的提升，写 \(\widetilde F=A(x,\theta)^{\mathsf T}+u(x,\theta)\)。\(u\) 关于 \(\theta\) 周期，且在两条边界上为零，所以 \(Du\) 在基本矩形上的积分为零。设该矩形面积为 \(A_1=2\pi L_1\)，由逐分量平方积分得
	\[
		\int_{C_1}|D\widetilde F|_{\mathrm F}^2\ge A_1|A|_{\mathrm F}^2
		=A_1(\alpha^2+\beta^2+1),
	\]
	这里 \(|\cdot|_{\mathrm F}\) 表示矩阵各项平方和的平方根。面积换元给出 \(\int_{C_1}J_F=2\pi L_2=\alpha A_1\)。若 \(K=K(F)\)，逐点的奇异值比不超过 \(K\)，故
	\[
		|D\widetilde F|_{\mathrm F}^2\le(K+K^{-1})J_F.
	\]
	积分并比较得到 \(K+K^{-1}\ge S\)。函数 \(k+k^{-1}\) 在 \([1,\infty)\) 上严格递增，因而 \(K\ge K_0\)，仿射映射达到下界。

	虽然 \(T\) 与 \(T+2\pi m\) 给出相同的边界点映射，它们指定了不同的提升，跨越圆柱的弧会多绕 \(m\) 圈。相对边界的同位保留这一次数，所以本题的标记类仍能记录整数扭转。取 \(T=0\)，公式回到无扭转环域的 \(\max(\alpha,\alpha^{-1})\)。
\end{solution}

\begin{exercise}
	固定 \(\lambda\ne0,1\)，求球面上在 \(0,1,\lambda,\infty\) 至多有简单极点、其他地方全纯的二次微分空间，并计算其非零元素在有限三个极点处的主部。
\end{exercise}
\begin{solution}
	写二次微分为 \(q=f(z)\,\mathrm{d}z^2\)。有限处的条件使 \(P(z)=z(z-1)(z-\lambda)f(z)\) 为整函数。在无穷远用 \(u=1/z\)，至多简单极点意味着 \(f(1/u)u^{-4}=O(u^{-1})\)，即 \(f(z)=O(z^{-3})\)。故 \(P\) 有界，由 Liouville 定理为常数。因此该空间是一维，基可取
	\[
		q_\lambda=\frac{\mathrm{d}z^2}{z(z-1)(z-\lambda)}.
	\]
	在 \(0,1,\lambda\) 处，\(\mathrm{d}z^2/(z-a)\) 的系数分别为 \(1/\lambda\)、\(1/(1-\lambda)\)、\(1/(\lambda(\lambda-1))\)。乘同一个常数即得任意元素的主部。无穷远的 \(u\) 坐标系数为 \(1/[u(1-u)(1-\lambda u)]\)，也恰有简单极点。
\end{solution}

\begin{exercise}
	设 \(X\) 为 \(y^2=\prod_{j=1}^6(z-a_j)\) 的光滑紧化，\(a_j\) 两两不同。证明
	\[
		H^0(X,K_X^2)=\operatorname{span}\left\{
		\left(\frac{\mathrm{d}z}{y}\right)^2,
		z\left(\frac{\mathrm{d}z}{y}\right)^2,
		z^2\left(\frac{\mathrm{d}z}{y}\right)^2\right\}.
	\]
	求对合 \((z,y)\mapsto(z,-y)\) 在此空间及形变切空间上的作用。
\end{exercise}
\begin{solution}
	由\cref{g4:cor:0204:hyperelliptic-genus}，\(X\) 的亏格为2，故本章维数公式给出 \(h^0(K_X^2)=3\)。在有限分支点可取 \(z-a_j=t^2\)、\(y=t\,c(t)\)、\(c(0)\ne0\)，于是 \(\mathrm{d}z/y=2\,\mathrm{d}t/c(t)\) 正则。在两个无穷远点取 \(u=1/z\)，有 \(y=u^{-3}v(u)\)、\(v(0)\ne0\)，从而
	\[
		z^k(\mathrm{d}z/y)^2=u^{2-k}v(u)^{-2}\,\mathrm{d}u^2.
	\]
	\(k=0,1,2\) 时均正则。三者线性无关，因为在没有分支的开集上任何线性关系都会给出恒零的二次多项式。因此它们为基。对合将 \(\mathrm{d}z/y\) 变号，平方后不变，并固定 \(z\)，所以在二次微分空间上作用恒等。Serre 配对在自同构下不变，故在对偶的形变切空间上也恒等。非平凡自同构在形变切空间的表示因而可以不忠实。
\end{solution}

\begin{exercise}
	设一个紧连通定向曲面的亏格为 \(g\)，有 \(n\) 条编号的测地边界，且 \(2g-2+n>0\)。若将其剪成双曲裤子，求裤子块数和内部拼接曲线数；在全部外边界长度固定或允许独立变化时，分别计算长度、扭转参数的实数量。
\end{exercise}
\begin{solution}
	设裤子有 \(p\) 块、内部拼接曲线有 \(k\) 条。每块 Euler 示性数为 \(-1\)，沿圆周粘合不改变总示性数，故 \(-p=2-2g-n\)，即 \(p=2g-2+n\)。每条内部曲线在切开后计两次，每条外边界计一次，故 \(3p=2k+n\)，于是 \(k=3g-3+n\)。每条内部曲线贡献一个正长度和一个实扭转，因此固定外边界时有 \(2k=6g-6+2n\) 个实参数；允许外边界长度变化再加 \(n\)，得到 \(6g-6+3n\)。外边界没有另一块裤子与之拼接，故不为它另加扭转参数。这里不固定边界上的逐点参数化；若另加边界点的标记，等价关系与参数也须相应调整。
\end{solution}

\begin{exercise}
	在闭曲面的带标记裤子分解中，固定所有长度和除一条曲线外的所有扭转。该曲线长为 \(\ell\)，所取扭转为 \(\tau\)。比较参数 \(\tau+m\ell\)、\(m\in\mathbb Z\) 给出的无标记曲面与带标记曲面。
\end{exercise}
\begin{solution}
	边界的弧长坐标属于 \(\mathbb R/\ell\mathbb Z\)，所以 \(\tau\) 与 \(\tau+m\ell\) 给出完全相同的边界点粘合，因而无标记的复曲面相同。实提升的变化却记录了跨越该环带的标记弧多绕了 \(m\) 圈，即改变了相应 Dehn 扭转的幂。带标记参数恢复保留这个实提升，因此这些参数在 \(m\) 不同时代表不同的带标记点。这个例子说明，在模空间中比较一个序列是否离开紧集，不能只看某个带标记扭转坐标是否趋于无穷。
\end{solution}

\begin{exercise}
	对结点光滑化 \(xy=t\) 作基变换 \(t=s^m\)，其中 \(m\in\mathbb Z_{\ge1}\)。给定 \(c>0\)，只考虑满足 \(0<|s|^m<c^2\) 的非零参数。判定总空间 \(xy=s^m\) 在原点是否光滑，并计算 \(|x|,|y|<c\) 中非零纤维的环域模量及两端相对转角。
\end{exercise}
\begin{solution}
	定义函数为 \(F(x,y,s)=xy-s^m\)，其梯度为 \((y,x,-ms^{m-1})\)。\(m=1\) 时梯度在原点非零，总空间光滑。\(m\ge2\) 时梯度为零；还可直接排除光滑性：曲线 \((t,0,0)\)、\((0,t,0)\)、\((t,t^{m-1},t)\) 都位于总空间内，其原点速度张成 \(\mathbb C^3\)。若原点是复子流形的光滑点，其切空间必须包含这三个方向，因而是三维；该子流形便在 \(\mathbb C^3\) 中局部开，与非零方程 \(xy-s^m=0\) 矛盾。对 \(s\ne0\)，\(y=s^m/x\)，所取区域为 \(|s|^m/c<|x|<c\)，模量为
	\[
		M(s)=\frac1{2\pi}\log\frac{c^2}{|s|^m}.
	\]
	角度满足 \(\arg y=m\arg s-\arg x\)。参数 \(s\) 绕零一圈时，两端的相对角度累计变化 \(2\pi m\)。每个非零纤维仍光滑，中心纤维仍为两分支结点；总空间的奇性与纤维的奇性因而是不同的判断。
\end{solution}

\begin{exercise}
	令 \(J(\lambda)=256(1-\lambda+\lambda^2)^3/[\lambda^2(1-\lambda)^2]\)。证明它在重排四点时不变，并给出 \(J:\mathbb P^1\to\mathbb P^1\) 的次数、全部分支值及其纤维重数。
\end{exercise}
\begin{solution}
	代入 \(1-\lambda\) 与 \(1/\lambda\)，分子分母的幂次相消，\(J\) 不变；这两个变换生成所需六元群。分子分母互素，次数为 \(6\)。无穷远值的原像为 \(0,1,\infty\)，各为二阶极点；零值的原像是 \(\lambda^2-\lambda+1=0\) 的两个根，各重数3。恒等式
	\[
		J(\lambda)-1728=\frac{64(\lambda+1)^2(\lambda-2)^2(2\lambda-1)^2}{\lambda^2(1-\lambda)^2}
	\]
	给出1728的三个原像 \(-1,2,1/2\)，各重数2。这些点的总分歧为 \(3(2-1)+2(3-1)+3(2-1)=10\)，已等于六次球面映射的 Riemann--Hurwitz 总数 \(2\cdot6-2\)，所以没有别的分支点。一个一般重排轨道有六点，都属于同一纤维，而一般纤维也只有六点，故 \(J\) 正好给出一般无序四点对象的参数；特殊纤维的重数反映额外对称与三个边界的合并。
\end{solution}

\begin{exercise}
	设 \(X\) 为亏格 \(g\ge2\) 的紧黎曼曲面，\(q\) 为非零全纯二次微分。若其零点处水平叶状结构的叉数为 \(p_1,\ldots,p_r\)，证明
	\[
		\sum_{j=1}^r(p_j-2)=4g-4.
	\]
	推出奇数叉点的个数为偶数，并列出亏格二时这一次数关系允许的叉数组合。
\end{exercise}
\begin{solution}
	在 \(m\) 阶零点附近，\(q=z^m a(z)\,\mathrm{d}z^2\)、\(a(0)\ne0\)。选 \(a\) 的全纯平方根后积分 \(z^{m/2}\sqrt{a(z)}\)，局部自然坐标的首项为非零常数乘 \(z^{(m+2)/2}\)。取一个全纯且非零的根因子重新选 \(z\) 坐标，可化成同一幂函数，所以叉数为 \(p=m+2\)。全纯截面 \(q\) 的零除子次数等于 \(\deg K_X^2=4g-4\)，将各 \(m_j=p_j-2\) 相加即得公式。模2化简可见奇数 \(p_j\) 的个数为偶数。

	亏格二时，零点阶数的正整数分拆为 \(4\)、\(3+1\)、\(2+2\)、\(2+1+1\)、\(1+1+1+1\)，因此次数关系允许的叉数组合为
	\[
		(6),\quad(5,3),\quad(4,4),\quad(4,3,3),\quad(3,3,3,3).
	\]
	这是次数所给的必要条件。正文由二叶分歧覆盖得到的伪 Anosov 模型具有两个四叉点，正好满足 \((4-2)+(4-2)=4\)。
\end{solution}

\begin{exercise}
	在定向环面的整数同调中取辛基 \(a,b\)，\(a\cdot b=1\)，令 \(v=pa+qb\)、\(\gcd(p,q)=1\)。考虑剪切 \(A_v(x)=x+(v\cdot x)v\)。求其矩阵及幂，并证明在 \(H_1(S;\mathbb Z/N\mathbb Z)\) 上的作用阶数恰为 \(N\)，其中 \(N\ge2\)。
\end{exercise}
\begin{solution}
	写 \(x=ra+sb\)，则 \(v\cdot x=ps-qr\)，所以
	\[
		A_v=I+B,\qquad B=\begin{pmatrix}-pq&p^2\\-q^2&pq\end{pmatrix}.
	\]
	因 \(v\cdot v=0\)，有 \(B^2=0\)，从而 \(A_v^k=I+kB\)，对正负整数均成立。若模 \(N\) 后 \(A_v^k=I\)，则 \(N\) 同时整除 \(kp^2,kpq,kq^2\)。由 \(p,q\) 互素，三个平方或乘积的最大公因子为1，它们的整数线性组合可为1，故 \(N\mid k\)。反之 \(N\mid k\) 时显然恒等。这说明整数同调上的无限阶扭转，在任一固定有限系数的同调中只能留下有限阶信息。若采用逆向扭转，矩阵换为 \(I-B\)，阶数结论不变。
\end{solution}
